\documentclass[10pt]{article}
\usepackage[T1]{fontenc}
\usepackage[utf8]{inputenc}
\usepackage[margin=1in]{geometry}
\usepackage{microtype}
\usepackage{amsmath,amssymb,amsthm}
\usepackage{booktabs,array,longtable}
\usepackage{xurl}
\usepackage[hidelinks]{hyperref}
\usepackage{graphicx}
\setlength{\emergencystretch}{2em}
\newtheorem{proposition}{Proposition}
\newtheorem{definition}{Definition}
\newcommand{\method}{Prometheus}
\newcommand{\R}{\mathbb{R}}
\newcommand{\E}{\mathbb{E}}
\title{Prometheus: Contextual Computational Primitives\\with Continuous State and Content Memory}
\author{Prometheus Research Project}
\date{Working manuscript: 2 November 2026}

% Canonical method: water/method.md. Direction: water/now.md.
% Protocol authority: ALL.md, TODO 5. Source and evidence remain versioned.
% Main text follows the human-specified eight-section paper structure.
% All detailed development records and unvalidated extensions are in the appendix.
\begin{document}
\maketitle
\begin{abstract}
How can a learning system form reusable computational operations and organize
their execution from task feedback? We investigate Prometheus, a trainable
Block combining continuous internal state, low-rank nonlinear residual
primitives, contextual operation calls, and content-addressed key--value
memory. Encoding, primitive functions, calling, memory writes and readout learn
jointly from terminal labels. Symbolic embeddings and learned continuous-input
adapters enter the same core. Soft execution combines primitive residuals;
hard-call execution selects one primitive while retaining continuous state and
memory. Their predictions and costs are evaluated separately. We give explicit
parameter, recurrent-state and forward-work accounts, and preserve earlier
topological protection constructions as optional mechanisms with distinct
assumptions. Fixed-budget development includes context-dependent sequence
tasks, sampled next-character prediction, and initial audio and image
classification diagnostics. In a parameter-matched 300-update character study,
the dynamic Block has perplexity 12.46 at context 32 and 12.50 at context 128,
compared with 12.80 and 48.41 for the tested Transformer and approximately 10.09
for the tested GRU. A same-size static-caller Block performs comparably to the
dynamic model. These single-seed observations motivate further mechanism and
resource comparisons; first real-input diagnostics favor the larger
Transformer controls. They establish finite implementations and measurements,
rather than broad competitive or continual-retention claims. A complete
Primary architecture remains to be frozen before the formal evaluation.
The central empirical question is whether this organization can improve the
acquisition--retention--resource tradeoff across meaningful task families.
\end{abstract}
\section{Introduction}
\label{sec:introduction}

How can a learning system construct, compose, and improve its own internal
computational mechanisms? Neural architectures already learn representations,
select information, implement memory, and organize reusable computation.
Attention, for example, learns input-dependent interactions through shared
projections and weighted aggregation~\cite{vaswani2017attention}. The question
is how computational capabilities can accumulate under continued learning:
which internal operations are acquired, how later tasks reuse them, and what
must be preserved when the parameters implementing those operations change.
An adequate answer should expose the executed computation and its learning
history, together with the resources required to retain it.

This problem has two coupled scales. At the implementation scale, a shared
parameter update changes representations, operations, and their invocation.
At the semantic scale, a capability is a complete input--computation--output
relation. Preserving an isolated operator is insufficient if its old inputs
are subsequently encoded differently or routed elsewhere. Conversely,
preserving all implementation details can prevent useful adaptation or require
separate capacity for each task. Dynamically expandable networks illustrate
that selective sharing, structural growth, and protection can be organized
jointly~\cite{yoon2018lifelong}. Their existence motivates a precise comparison
of learning organizations under shared information and resource conditions.

We investigate \emph{Prometheus}, an executable Block with a continuous internal
state, learned nonlinear residual primitives, a contextual caller, and
content-addressed memory. The caller reads the current state, input and memory
content to select or mix primitive updates. After those calls, learned
addresses and gates revise the memory's keys and values. A task supplies
observed inputs and a terminal label; gradients jointly train representation,
operations, invocation, memory and readout. No task-specific primitive table
or correct calling sequence is installed by the experimenter.

The architecture makes two interfaces explicit. First, symbolic embeddings
and modality-specific projections enter the same core, permitting direct tests
of a common computation on finite-state programs, text prefixes, audio and
image features. Second, soft mixture execution and hard-call execution use
the same parameters but follow distinct trajectories. Hard calls retain a
continuous state and readout. Their quality and costs require separate
measurement; replacing a mixture with its largest component is an empirical
change to the computation.

This continuous design follows measured limitations of earlier categorical
prototypes, including weak terminal credit and collapsed calling behavior.
Residual state updates and learned memory are concrete architectural responses,
whose effects can be tested against static calling, expanded state, recurrent
networks, Transformers and applicable state-space models. The component
operations are established mechanisms. The scientific target is an important
acquisition--retention--resource benefit from their organization, with clear
causal attribution. Historical monodromy-based protection remains an optional
research mechanism documented with its original assumptions and controls.

Current evidence consists of bounded development runs. A contextual sequence
experiment demonstrates learning with both execution interfaces; sampled
character prediction tests length changes and parameter-matched controls;
initial audio and image diagnostics exercise the continuous-input interface.
The results are mixed and have not yet established continuing capability
growth or a broad competitive advantage. This manuscript therefore connects
an implemented, inspectable architecture to a prospective continual-learning
evaluation. A validated Primary freeze and subsequent retraining separate
exploratory changes from the eventual empirical claims.

\section{Related Work}
\label{sec:related-work}

\paragraph{Context-dependent neural computation.}
Transformers learn input-dependent interactions through attention and combine
them with shared feedforward transformations~\cite{vaswani2017attention}.
Prometheus also uses task-driven neural representations and weighted content
selection. Its caller chooses among low-rank nonlinear residual functions;
its externalized key--value state supplies a learned read vector to those
functions. The static-caller control retains the same contextual primitives
and memory, separating the effect of context-dependent invocation from the
presence of a recurrent memory system. This comparison can identify an effect
of the caller without treating memory, function composition or expert choice
as new capabilities in themselves.

Recurrent and selective state-space models provide additional conventional
organizations of context. Gu and Dao's selective SSM makes step sizes and
write/read vectors depend on the current input, controlling how information
enters and propagates through state.\footnote{Original mechanism: Algorithm2
and Sections3.2--3.4 of \url{https://arxiv.org/abs/2312.00752}.}
Its state parameterization and recurrence differ from the explicit primitive
mixture considered here. Comparisons must also distinguish an author's
optimized implementation from a reference adaptation: mathematical recurrence
alone does not reproduce fused scan, recomputation, or kernel performance.

\paragraph{Continual organization.}
DEN selectively retrains relevant subnetworks, conditionally expands capacity,
and uses splitting and timestamped inference to manage semantic
drift~\cite{yoon2018lifelong}. Replay, freezing and modular capacity are likewise
concrete comparison conditions in our protocol. The present continuous Block
requires its own retention measurements during shared learning. Parameters,
stored examples, auxiliary states, maintenance and execution must all enter
comparisons; equivalent terminal supervision does not automatically equate
total work or attained capability.

\paragraph{Discrete mechanisms and topological protection.}
Earlier Prometheus constructions used class-boundary feedback for loop
operations. Their derivative decomposition follows established stochastic
computation-graph and discontinuous-gradient
analyses~\cite{schulman2015stochastic,lee2018reparameterization}; it is distinct
from ordinary differentiable training of the current residual functions.
Monodromy supplies a conditional operation-invariance mechanism, retained in
the appendix as an optional component. Kitaev's braiding construction concerns
physical topological quantum protection~\cite{kitaev2003fault}; classical loop
simulation does not inherit that hardware guarantee. Bennett's reversible
simulation explicitly retains and later cleans auxiliary
information~\cite{bennett1973logical}. It informs historical workspace analyses,
whereas the current nonlinear state and memory overwrite rules directly
allow irreversible updates. Each organization needs its own learning,
retention and resource evidence.

\section{Prometheus Architecture}
\label{sec:architecture-main}

Prometheus investigates a trainable computational Block whose internal
operations, calls, and memory are learned together from task feedback. The
current development architecture uses a continuous state, a bank of nonlinear
residual primitives, and content-addressed memory. This revises the earlier
categorical and loop-based implementations without transferring their
protection guarantees to the new model. Its ingredients are established neural
mechanisms; the question is whether their particular organization supports
useful acquisition, composition, and retention at competitive total cost.

\subsection{State, inputs, and executable primitives}

Let $h_{t,s}\in\mathbb R^H$ denote the state before call slot $s$ of input $t$.
Memory contains keys $\mathsf K_t\in\mathbb R^{M\times Q}$ and values
$\mathsf V_t\in\mathbb R^{M\times H}$. An observed initial-state identifier
selects a learned embedding $h_{1,0}=E_\theta[x]$. Keys begin at a learned,
task-independent matrix and values at zero for each example. There is no
persistent memory across examples in this implementation.

The raw input feature $u_t\in\mathbb R^T$ is a token embedding for symbolic
inputs or the output of a learned modality adapter. A learned slot embedding
gives $x_{t,s}=u_t+e_s$. All slots reuse the same primitive bank and caller.
During a token, memory is fixed while the state changes. Each call computes
\begin{align}
 r_{t,s}&=\operatorname{softmax}\!\left(
   q_\theta(h_{t,s},x_{t,s})\mathsf K_t^\top/\sqrt Q\right)\mathsf V_t,
   \\[-1mm]
 \alpha_{t,s}&=\operatorname{softmax}
   f_\theta(h_{t,s},x_{t,s},r_{t,s}),\\
 z_{t,s}&=[h_{t,s};x_{t,s};r_{t,s}],\\
 \delta_{t,s,k}&=U_k\operatorname{GELU}(D_k z_{t,s}+b_k),\\
 h_{t,s+1}&=\operatorname{LayerNorm}\!\left(
    h_{t,s}+\sum_{k=1}^{K}\alpha_{t,s,k}\delta_{t,s,k}\right).
 \label{eq:main-block-computation}
\end{align}
Here $D_k\in\mathbb R^{R\times(2H+T)}$ and
$U_k\in\mathbb R^{H\times R}$. Each primitive is an executable nonlinear
state update with its own learned parameters. The bank has no supplied
task-specific operation or explicit identity choice. A residual path replaces
repeated row-stochastic state propagation; neither normalization nor low rank
implies invertibility or exact preservation of information.

\subsection{Content memory and composition}

After the last call, write features concatenate the final state, raw input,
and last read vector: $\bar z_t=[h_{t,S};u_t;r_{t,S-1}]$. Learned projections
produce an address, gate, candidate key and candidate value:
\begin{align}
 a_t&=\operatorname{softmax}
    (w_\theta(\bar z_t)\mathsf K_t^\top/\sqrt Q),
       &g_t&=\sigma(g_\theta(\bar z_t)),\\
 \widehat k_t&=\tanh(k_\theta(\bar z_t)),
       &\widehat v_t&=v_\theta(\bar z_t).
\end{align}
With $\beta_{tj}=g_ta_{tj}$, each slot updates as
\begin{equation}
 (\mathsf K_{t+1,j},\mathsf V_{t+1,j})=
 (1-\beta_{tj})(\mathsf K_{t,j},\mathsf V_{t,j})+
 \beta_{tj}(\widehat k_t,\widehat v_t).
\end{equation}
The next token starts from $h_{t+1,0}=h_{t,S}$. Thus later calls depend on
the actual intermediate computation and retained content. Order-sensitive
composition follows from executing successive nonlinear updates; no
commutativity constraint is imposed. The terminal distribution is
$\operatorname{softmax}(W_Rh_{L,S}+b_R)$.

The development configuration has $H=64$, $T=Q=16$, $M=8$, $K=4$, $R=8$ and
one call per token. All dimensions are explicit resources, rather than counts
inferred from a hidden task generator. A static-caller control zeros only the
state and memory-read inputs to $f_\theta$; its primitives, memory, shapes and
parameter count are unchanged. It therefore retains contextual computation
and is stronger than a memoryless ablation.

\subsection{A common input interface}

Symbolic lookup and externally embedded inputs enter the same execution
function. Audio features or image patches require a learned projection to
$\mathbb R^T$, whose parameters and work are charged separately. No modality
adapter supplies operations, calls, intermediate targets, or task identities.
This common interface permits testing the same core across tasks; it does not
by itself establish transfer or continual retention. The complete Primary
architecture remains unfrozen, and present measurements are separately
identified development evidence.

\section{Learning and Retention}
\label{sec:learning-main}
\label{sec:monodromy-learning-main}

\subsection{Joint operation formation from terminal feedback}

For observed examples $(x,u_{1:L},y)$, the current learner minimizes terminal
negative log likelihood,
\begin{equation}
 \mathcal L(\theta)=-\frac1B\sum_{i=1}^{B}
    \log p_\theta(y_i\mid x_i,u_{i,1:L_i}).
\end{equation}
The initial representation, primitive matrices, caller, memory read/write
projections and readout all receive gradients through the executed sequence.
The learner receives neither correct calls nor intermediate states. Primitive
formation here means changes to executable learned state updates, rather than
assignment of examples to an already supplied answer table. It does not imply
that a primitive becomes a separately named or independently interpretable
algorithm. Such behavior requires execution and intervention evidence.

Soft execution evaluates every primitive and combines its residual according
to the learned call probabilities. This supplies terminal credit to the whole
bank. Backpropagation through time also differentiates the learned content
memory; state and memory reset for each example. The continuous development
runs use soft training. Their optimization budget, sampled minibatches and
validation protocol are fixed independently of test scores. Learning curves
and failed predecessors are retained rather than replacing low results with
an assumed successful configuration.

\subsection{Soft and hard-call execution}

Hard-call execution selects
$\widehat k_{t,s}=\arg\max_k\alpha_{t,s,k}$ and evaluates only
$\delta_{t,s,\widehat k_{t,s}}$. Encoding, state, memory, primitive functions
and the probabilistic readout remain continuous. We call this interface
\emph{native} to distinguish it from mixture execution; it is not a fully
discrete program or a one-hot classifier. Both interfaces use the same saved
weights, and both yield normalized terminal distributions.

An optional straight-through interface uses hard calls in the forward pass
and a dense soft-mixture surrogate in the backward pass. Writing
$\delta^{\mathrm{soft}}=\sum_k\alpha_k\delta_k$ and using
$\operatorname{sg}$ for stop-gradient, its residual is
\begin{equation}
 \delta^{\mathrm{ST}}=
 \operatorname{sg}(\delta_{\widehat k})+
 \delta^{\mathrm{soft}}-\operatorname{sg}(\delta^{\mathrm{soft}}).
\end{equation}
This is a biased surrogate for a discontinuous decision, not its exact
gradient. Soft and native trajectories can diverge because subsequent calls
and memory writes depend on their resulting states. Consequently, likelihood,
accuracy and inference cost must be measured separately; correctness of soft
execution does not certify a hard-call program.

\subsection{Retention as an empirical obligation}

Shared updates may change old encodings, primitive behavior, calls, memory
access and readout. The present continuous Block has no built-in theorem
protecting their composition. Its continual-learning policy must therefore
be evaluated against replay, freezing, modular adaptation and other applicable
controls with the same information. Extra examples, snapshots, auxiliary
state and maintenance are counted. Finite current-task success cannot replace
measurements during and after later learning.

Earlier implementations supply a distinct optional protection mechanism:
homotopic deformations of a based loop preserve its monodromy action, while
additional finite closure checks protect learned task bindings. Those
constructions and their counterexamples remain in
Appendices~\ref{sec:monodromy-principle} and~\ref{sec:prometheus-block}.
Their certificates apply to the specified finite execution models, not to the
unconstrained nonlinear primitives above. Monodromy is not required for the
current Primary candidate; its later inclusion would need an explicit
integration and a measured benefit beyond strong ordinary controls. Neither
this optional component nor an independent new learning rule is presumed
necessary for a useful architecture.

\section{Theoretical Analysis}
\label{sec:theory-main}

We analyze the implemented resource structure, separating trainable weights,
transient state and execution. These are conditional accounting identities
and explicit work models, not claims of superior complexity or protection.

\paragraph{Parameters and working state.}
Let $N$ be the output/initial-state cardinality, $V$ the input vocabulary,
$H$ the state width, $T$ the token width, $Q$ the key width, $M$ memory slots,
$K$ primitives, $R$ primitive rank, $C$ caller width and $S$ calls per token.
Set $F=2H+T$. Including biases, the current Block has
\begin{align}
 P={}&NH+VT+ST+MQ+\big((H+T)Q+Q\big)\nonumber\\
    &+(FC+C+CK+K)+K(RF+R+HR)+2H\nonumber\\
    &+2(FQ+Q)+(FH+H)+(F+1)+(HN+N).
 \label{eq:main-continuous-parameters}
\end{align}
The terms account respectively for learned interfaces and initial keys,
queries, calls, primitives, normalization, writes and readout. A modality
projection adds its own parameters. Increasing sequence length adds no model
weights, while increasing $H$, $M$, $K$ or rank changes both capacity and cost.

The recurrent payload for one example is
\begin{equation}
 C_{\mathrm{state}}=H+M(Q+H)
\end{equation}
scalars. For the stated development dimensions this is 704 floats. It excludes
temporary tensors, gradients, optimizer moments and saved training activations.
A streaming implementation can retain this state and process a new input;
the current forward API instead resets per example. Current character
experiments recompute each entire observed prefix. Their timings therefore
do not establish a streaming or autoregressive caching advantage.

\paragraph{Forward work and training memory.}
For $B$ examples and maximum supplied length $L$, dense primitive projections
contribute at most
\begin{equation}
 W_{\mathrm{primitive}}=BLSKR(F+H)
\end{equation}
multiply--accumulate operations in soft execution, versus
$BLSR(F+H)$ for actual selected-expert evaluation. This conditional saving
does not equate their predictions. Memory reads additionally require query
projection and $O(BLSM(Q+H))$ attention products; writes and the caller have
their own projection and update terms. Source counters report these components
separately. They are forward MAC proxies, not measured training FLOPs: they
exclude normalization, activations, indexing, backward and backend workspace.
Actual wall time and allocator peaks remain necessary measurements.

Backpropagation through a prefix retains intermediate states and operator
activations. Its memory consequently grows with unrolled length even though
the inference carry is fixed. Parameter storage, recurrent payload, training
memory, checkpoint bytes and inference latency are distinct resources.

\paragraph{Composition and protection boundaries.}
Repeated calls define an explicit recurrent composition. Low-rank residual
updates, LayerNorm and content addressing do not guarantee exact inversion,
arbitrary memory recovery, length extrapolation, or preservation after shared
learning. Likewise, the finite stochastic-path certificate of the earlier
categorical model is not a certificate for this context-dependent continuous
model. Historical topological preservation and reversible-workspace results
retain their original assumptions in the appendix. The current architecture
must establish useful acquisition and retention through the matched empirical
comparisons of Section~\ref{sec:experiments}.

\section{Experiments}

The character-task control--value development Block reduces the model
from 41,593 to 20,721 parameters under the same 300-update all-token
protocol. Its paired 64/128/256-context perplexities are
10.593/10.498/10.560, versus 10.358/10.251/10.253 for the earlier Block.
Training-forward work decreases from 21.715 to 10.342 billion MAC proxies,
but measured training wall time changes from 75.96 to 79.29 seconds.
The GRU and official Mamba2 still obtain lower perplexity and lower
recorded training time. These distinct model versions are development
comparisons, not pooled Primary results or a complete-cost victory.
\label{sec:experiments}


A later three-arm memory-interface development test completed 7,200 new
updates. The 10,351-parameter joint repair achieved 62.70\% new-task IID,
but only 5 of 192 label-changing query pairs were correct at both endpoints.
Its 26.916\,s full-stream update time exceeded the original GRU32's
17.281\,s despite fewer parameters. A frozen VALID content-direction
intervention with exactly matched write weights and coverage reduced
62.50\% accuracy to 41.99\%; strength-matched cached-direction
compensation obtained 39.06\%. These measurements motivate continued
binding research, not a Primary acceptance or comparative efficiency claim;
complete protocols and costs are retained in Appendix~\ref{sec:prometheus-block}.

Our evaluation asks whether contextual computational primitives improve
acquisition, retention, or resource use relative to strong conventional
architectures. Topological components receive a separate causal test when
included. The primary experimental object is the latest fully fused
Prometheus architecture after general Block validation. Monodromy is an
optional component selected by measured benefit; its independent advantage
is not a prerequisite for the Primary freeze. Its core mechanism,
code version, and evaluation configuration will be frozen before retraining
across all five evaluation layers. No publication results for the to-be-frozen
Primary are available at the present manuscript checkpoint. Earlier coordinate-renderer
streams and small integrated diagnostics are development evidence in
Appendices~\ref{app:benchmarking} and~\ref{app:evidence}; they are not pooled
with Primary results or used to supply its performance or efficiency claims.

\subsection{Observed development results}

We first tested the same continuous core on three-stage contextual programs,
sampled character prediction, and initial real-input recognition. These are
fixed-budget, single-seed development measurements, not formal Primary results.
The context stream exposes raw tokens and an initial state, with terminal labels
but no task identifier, correct route, or intermediate state. Training lengths
are 2--12; separate tests use lengths 24--48 and structured unseen compositions.
Each phase uses 600 updates of 64 examples. Actual phase-0 weights, Adam state,
and RNG were restored for later phases; phase0 was not retrained.

\begin{table}[t]
\centering\small
\begin{tabular}{lrrrrr}
\toprule
Model & Parameters & Task0 & Task1 & Task2 & Train s\\
\midrule
Continuous dynamic &28,365&82.81&90.23&91.99&92.21\\
Continuous static caller &28,365&81.45&94.92&94.14&79.12\\
GRU44 &25,572&93.95&95.90&93.36&10.25\\
Transformer32 &26,728&48.63&79.30&83.01&13.76\\
\bottomrule
\end{tabular}
\caption{Final IID accuracy (\%) after three development phases, without
replay. Train s sums actual gradient-step wall time, not total work or exclusive
throughput. The static caller retains contextual primitives and memory.}
\label{tab:development-context}
\end{table}

The dynamic model forgets part of task0 and does not uniquely outperform the
static caller. Distribution checks are essential: a TRAIN-majority constant
predictor scores 69.73\% on task2's long test, equal to the dynamic model's
69.73\%. Its structured-unseen scores are 25.78\% and 25.98\%, respectively.
Thus high IID accuracy alone does not establish general memory computation.

All-token character development uses 300 updates of 32 length64 windows,
614,400 next-character labels per method. A causal-mask Transformer and GRU
receive identical windows and batches. Future-token perturbations test actual
causality before training. The 512 frozen TEST endpoints are paired across
lengths64/128/256; length comparisons score the same final 64 labels. These are
sampled-window character PPLs, not full-corpus or word-level PPLs.

\begin{table}[t]
\centering\small
\begin{tabular}{lrrrrr}
\toprule
Model &P&PPL64&Train s&Driver s&Forward G MAC\\
\midrule
Dynamic &41,593&10.358&75.96&119.47&21.715\\
Static &41,593&10.306&75.24&118.74&21.715\\
Detached &41,593&10.243&74.07&114.41&21.715\\
Transformer36 &39,664&10.533&6.75&12.26&29.816\\
GRU44 &32,660&8.664&2.49&7.96&16.977\\
Official Mamba2 &41,116&7.173&3.56&4.58&---\\
Control--value Block &20,721&10.593&79.29&117.67&10.342\\
Same Block, training Graph &20,721&10.593&7.62&47.72&10.342\\
GRU34 &21,180&9.449&1.09&5.48&10.612\\
Transformer24 &19,564&11.414&3.12&7.78&15.630\\
\bottomrule
\end{tabular}
\caption{Actual all-token language development. Driver time includes measured
evaluation/checkpoint work but excludes process startup. MAC proxies count
training forwards only, not backward or total FLOPs; all raw costs are retained.}
\label{tab:development-ar-cost}
\end{table}

The two smaller conventional widths were fixed prospectively from parameter
formulas, with no accuracy gate or TEST-based configuration selection.
They use the original windows, feedback, Adam, clipping and 300-update budget.
This parameter-matched comparison narrows the apparent weight advantage:
the 20,721-parameter Block remains worse than GRU34 in paired perplexity and
recorded training time. All widths remain separate development runs;
they are not pooled as seeds or substituted for a frozen Primary.
The later Graph run retains the exact core, data, labels and 300 batch indices.
Actual forward/backward replay with eager Adam reduces update wall time by
10.40 times, while preserving paired perplexity to the reported precision.
Setup, diagnostic gradients, memory pools and evaluation are charged separately:
32.16 seconds of evaluation still dominates its 47.72-second driver.
This is an execution improvement, not a new learning rule or a baseline victory.

With contexts of 64, 128, and 256 characters, dynamic common-suffix PPL is
10.358, 10.251, and 10.253; Transformer yields 10.533, 15.185, and 14.736,
and GRU yields 8.664, 8.552, and 8.552. Stability under longer
inputs alone does not prove useful distant-memory use. The earlier terminal-prefix
protocol and serial eager/hard-call inference audits remain separately reported
in the appendix and are not pooled into this new protocol. They show no
hard-call speedup. Cached streaming serving has not been established.
The subsequently executed official Mamba2 baseline uses the same 300-update
protocol and achieves common-suffix PPL 7.173/7.048/7.049 with 41,116 parameters.
Its 3.56-second training and 4.58-second driver times exclude 82.67 seconds of
measured cold compilation/priming; the entire coordinator launch takes 93.31
seconds and creates 55.19 MB of kernel cache. Its exact selective-SSD and
backward FLOPs remain unmeasured. Earlier arms have their own recorded
concurrency, so this is not an isolated-throughput ratio. These results
provide no broad competitive advantage for the current Block.

\begin{table}[t]
\centering\small
\begin{tabular}{lrrrr}
\toprule
Model &Speech acc.&Image acc.&Speech/image P&Train s (S/I)\\
\midrule
Dynamic &29.77&42.10&32,280/40,793&97.51/66.28\\
Static caller &71.17&40.10&32,280/40,793&98.49/66.05\\
Detached &69.38&39.00&32,280/40,793&97.84/65.97\\
GRU64 &76.25&44.60&54,819/59,556&1.60/1.92\\
Transformer64 &76.48&49.90&105,379/110,116&7.92/7.85\\
CNN &59.77&44.30&35,683/38,852&.96/1.05\\
\bottomrule
\end{tabular}
\caption{Initial Speech Commands/CIFAR-100 development TEST accuracy (\%).
Each uses the first 10 official classes, 128 TRAIN examples per class, common
features and 300 updates. Full 35/100-class heads are charged. This is not yet
full-dataset or class-incremental validation. Training wall times reflect the
original recorded possible concurrency, not isolated throughput.}
\label{tab:development-real}
\end{table}

The detached variant preserves the dynamic forward rule and parameter count,
but stops only the caller-input gradients from state and memory. CPU autograd
checks verify this intervention. It improves the poor speech optimization
trajectory, while not improving the image result; it is not a general solution.
Training and evaluation work, optimizer state, input caches, and failed
versions are recorded separately. The current Python-unrolled Block is slower
than the competitors despite smaller forward MAC proxies. Consequently these
observations do not establish a broad capability or complete-resource advantage.

\subsection{Questions and comparison conditions}

We distinguish three questions: whether new computations are acquired from
terminal feedback, whether later learning preserves their complete execution,
and whether topology independently changes the acquisition--retention--resource
tradeoff. Preservation includes input encoding, operation selection, composition,
and readout, rather than operation maps alone. All methods receive the same
permitted examples, terminal labels, stage order, and inference information.
Teacher transitions, intermediate execution states, correct routes, and
preconstructed task-specific primitives are excluded as free supervision.
Discovery, installation, observation-table passes, and protection are charged.

We separately compare matched information, matched resource budgets, and matched
attained quality. Equal update counts alone do not equate total learning work
when a method also performs program discovery or verification. Accuracy--parameter,
accuracy--training-cost, and accuracy--inference-cost frontiers expose these
differences. Changes to input adapters and output heads are recorded and charged;
a changed core learning rule constitutes a different architecture variant.

\subsection{Tasks and continual streams}

\begin{table}[t]
\centering
\small
\begin{tabular}{p{0.09\linewidth}p{0.40\linewidth}p{0.40\linewidth}}
\toprule
Layer & Task family & Principal measurement\\
\midrule
I & Six-stage streams at $N=8,16,32$ & Acquisition, stage-boundary and
within-stage retention, parameter growth\\
II & Noncommuting programs, long compositions, key--value memory & Order
sensitivity, unseen composition, delayed state use\\
III & Reversible, nonreversible, continuous, variable-length computation &
Generality of the same core Block\\
IV & Speech Commands v2 and Split CIFAR-100 & No-task-ID class-incremental
recognition and forgetting\\
V & Architecture competitors and component interventions & Competitive
frontiers and independent mechanism attribution\\
\bottomrule
\end{tabular}
\caption{Five required empirical layers. Success on one does not replace
measurements at another. The complete protocol is in Appendix~\ref{app:benchmarking}.}
\label{tab:main-evaluation-layers}
\end{table}

The finite computational streams expose observable initial states, token
programs, public task context, and terminal outputs. Context-dependent token
semantics make that public context available to every competitor. We disclose
short-word coverage, repeated examples, and every additional label access.
Both differentiable neural execution and native deployment are evaluated on
new and old tasks. Intermediate old-task minima and their subsequent recovery
are reported separately from stage-boundary retention.

The composition suite includes $AB\ne BA$, commutators, and unseen operation
combinations. Training lengths are at most eight; tests at lengths 16, 32, and 64
are reported separately. Key--value tasks require writing, overwriting,
reading, and delayed queries from terminal supervision. Generality tests cover
merging or erasing operations, continuous inputs/outputs, and variable-length
execution. Reversible lifting, extra workspace, or external memory counts
toward resources rather than making irreversibility free.

The following formal protocol remains prospective. Speech Commands v2 uses 35 classes in seven phases of five, the official
train/validation/test split, and common 16-kHz preprocessing. Split CIFAR-100
uses 100 fine classes in ten phases of ten, the official 50,000 training and
10,000 test images, and a fixed class-stratified 10\% training holdout for
validation. All methods share preprocessing and augmentation. Both tasks
recognize all seen classes without a task ID or class-group identifier at
inference. Their Memory-0 and Memory-500 conditions permit zero or at most 500
stored old examples in total. Stored features, teacher outputs, and statistics
are additionally disclosed. Main comparisons start from random initialization;
pretraining is supplemental with equal access across methods.

\subsection{Competitors and bounded development}

The primary competing family is the Transformer: vanilla, experience replay,
adapters, and replay plus adapters, at two prospectively selected scales.
Continual-learning comparisons include DEN, PackNet, replay, and ER-ACE or a
mature comparable class-incremental algorithm. Task-incremental PackNet is
separated from no-ID class-incremental rankings. Mamba-2 participates in
sequence/memory tests, Titans in applicable memory tasks, and HOPE/Nested
Learning where a verifiable implementation supports the interface. Ordinary
program learning, freezing, modular expansion, and constraint optimization
provide additional matched-information controls; they do not replace neural
architecture comparisons.

\begin{table}[t]
\centering
\small
\begin{tabular}{llll}
\toprule
Setting & Low budget & Medium budget & High budget\\
\midrule
Synthetic, updates/phase & B1:600 & B2:2,400 & B3:9,600\\
Real data, epochs/phase &10&50&200\\
\bottomrule
\end{tabular}
\caption{Prospective budget tiers available to every method. Discovery,
maintenance, and configuration-development work are additional reported costs.}
\label{tab:main-training-budgets}
\end{table}

Each main family has at most 24 development configurations. Development seed 0
selects candidates using frozen validation criteria; the top three are checked
with development seed 1. Configuration lock precedes final seeds 0, 1, 2, 3, 4.
Test and OOD results never select hyperparameters, retries, stopping points,
or configurations. Methods may use their appropriate optimizers and schedules.
Extra-budget trajectories remain separately labeled supplemental configurations.
Only localized, reproducible software, data, or training-logic defects qualify
as bug fixes; ordinary tuning is development.

Configuration validity depends on checked implementation, fidelity of any
claimed reproduction, information fairness, and compliance with the prospective
protocol, never on an accuracy threshold. High, intermediate, and random-level
scores are all valid outcomes under those conditions. A low-scoring adaptation
supports conclusions about that adaptation and budget, not an entire published
family. Unfaithful reproductions and untested implementations are explicitly
identified; they are not made invalid by low performance. Failed acquisition,
all failed seeds, and additional resource use are retained without extending
trials until a favorable result appears.

Learning curves distinguish failure to acquire a first task from forgetting
after later tasks. Training and validation diagnostics examine representation
limits, optimization, and credit assignment without using held-out tests to
choose a repair. A substantive architecture or optimizer revision receives a
new configuration identifier and preserves its predecessor's results. Formation
before gradient training, maintenance after training, and admission are separate
events, with their own label accesses and costs. This prevents installed
programs or confidence sharpening from being misreported as capabilities first
discovered by the gradient updates.

\subsection{Attribution, costs, and statistical analysis}

Same-architecture interventions remove reuse, sharing, state growth,
old-capability protection, graph constraints, and native compilation separately.
Topology-specific controls substitute ordinary constraints, disable cross-class
loop learning, disable within-class safe deformation, and compare ordinary
boundary optimization or discrete search. Information, certificates, and caches
are shared when applicable. Benefits from compression, replay, ordinary closure,
or confidence maintenance are not attributed to monodromy. A topology benefit
requires a measured residual effect under these concrete controls.

We record total, trainable, and frozen parameters; actual model, optimizer,
cache, replay, reference, and checkpoint storage; updates and sample accesses;
training FLOPs and GPU time; discovery/search and rejected proposals; growth,
maintenance, protection, verification, and reload. Controlled same-hardware
inference reports latency p50/p95, throughput, and peak memory. Concurrent wall
times are not ranked as isolated throughput, and overlapping timing/storage
fields are not summed twice. Parameters per newly retained capability include
the initial allocation and do not hide costs in nonparameter state.

Final accuracy, average incremental accuracy, forgetting, backward transfer,
acquisition cost, and equal-quality resource use are reported over all five
independent final seeds. We report means, standard deviations, 95\% confidence
intervals, and paired-seed differences; multiple primary tests receive
multiplicity correction. Resumed fragments of one lineage are not independent
replicates. Until these frozen-Primary measurements exist, this protocol
specifies the evaluation rather than asserting an experimental advantage.

\section{Discussion \& Limitations}
\label{sec:discussion}

The current Block makes computational operations and their invocation explicit
within a trainable continuous system. Its contextual sequence development run
shows that terminal supervision can train useful nonlinear state updates and
calling behavior. The same-size static-caller model is also strong. Because
that control still gives the primitives and memory full contextual inputs,
these observations do not isolate a benefit of dynamic calling. The useful
next mechanistic question is which learned primitive functions, memory reads
and invocation decisions account for the acquired computation, and whether
those roles remain reusable when task demands change.

Parameter-matched character development makes the comparison more specific.
After the same 300 updates, the 41,593-parameter dynamic Block has sampled
perplexities 12.458, 12.468 and12.496 at contexts 32, 64 and 128, respectively.
The static variant has 12.383, 12.452 and12.436. The 43,300-parameter Transformer
has 12.803, 49.837 and48.409, while the 41,548-parameter GRU has approximately
10.090 across these contexts. Context 64 is the designated unseen-length split
and context 128 the long split. These are 512 paired sampled positions per split,
not whole-corpus or word-level perplexities. They show a stable contextual
behavior for the tested Block but also a stronger recurrent control. The
Transformer's long-prefix result characterizes this configuration and budget,
not the whole architecture family. Matching parameters does not match
training work, forward cost or transient memory.

Initial real-input diagnostics expose a further gap. With the first ten
official classes, at most 128 training examples per class and 300 updates,
the speech Block has approximately 29.77\% soft and 22.19\% hard-call accuracy,
compared with 76.48\% for the larger Transformer. The image Block obtains
approximately 42.1\% and 30.4\%, respectively, compared with 49.9\% for the
Transformer. Their Block parameter counts are 32,280 and 40,793; the Transformer
counts are 105,379 and 110,116. These static development tests exercise the
same core through learned adapters. They do not yet test full-class continual
recognition, and the soft/hard gaps show that inference selection remains a
material architectural choice.

The central remaining obligation is continual acquisition with complete
retention and resource accounting. The continuous Block currently resets
per-example memory, recomputes observed prefixes and has no general shared-update
protection certificate. Earlier finite monodromy and binding-closure results
retain their narrower scope in the appendix. Any integration must explain
what operation or task relation is protected in the new state space and
measure its benefit against ordinary constraints. An independent new primitive
is not a prerequisite for an important empirical result; a broad reproducible
advantage still requires a frozen architecture, strong appropriate controls,
and the full evaluation rather than a favorable subset of development scores.

\section{Conclusion}
\label{sec:conclusion}

Prometheus investigates how learned operations, contextual invocation and
content memory can organize executed computation. The current implementation
uses continuous state and low-rank nonlinear residual primitives, jointly
trained with representation, memory and readout from terminal labels. A
shared core accepts symbolic or continuous features and supports separately
measured soft and hard-call interfaces. Bounded sequence, sampled-character,
audio and image development provides concrete learning and resource evidence,
with mixed outcomes against recurrent and Transformer controls. It does not
yet establish broad advantage or protected continual capability growth.
Historical topological constructions and their negative controls remain
available as distinct optional mechanisms. The next scientific claim must
come from the complete architecture's acquisition, retention and cost under
frozen fair comparisons. The present manuscript records that architecture
and evidence boundary while the Primary design and formal evaluation remain
in progress.


\clearpage
\bibliographystyle{plain}
\bibliography{references}
\clearpage
\appendix
% Detailed task-corresponding materials are intentionally outside the main text.
\section{Continual acquisition and capability growth}
\label{sec:capability-growth}

\paragraph{Historical mechanism and Primary architecture boundary.}
This section records the earlier native/coordinate-renderer prototype as
historical mechanism evidence and internal ablations. It is not the Primary
Prometheus architecture for the publication benchmark. Primary Prometheus must
be the latest fully fused architecture actually validated under TODO~3; its
core mechanism, code version, and evaluation configuration must be frozen
before all five benchmark layers are retrained. Historical streams cannot be
pooled with those new Primary results or used to supply their parameter and
cost claims. The limited D3 hybrid V2 diagnostic is subsequent development,
not completed TODO~3 validation or a frozen Primary architecture.

The historical acquisition system learns finite recurrent computations from terminal
examples, installs the inferred operations in neural parameters, and subsequently
trains and maintains a differentiable calling interface. This is a hybrid
identification--installation--gradient-learning procedure. Its successful
synthetic executions establish a working accumulation mechanism within the
tested task family; they do not yet establish that topology improves acquisition
or parameter growth. In particular, an ordinary constrained implementation with
the same formation and maintenance rules matches the measured results.

\subsection{Historical representations, operations, and calling}

A task example consists of an observable initial state $x\in\{1,\ldots,N\}$,
a token word $w=(a_1,\ldots,a_L)$, a public context $c$, and a terminal label
$y\in\{1,\ldots,N\}$. The public context is necessary in these synthetic
streams because identical tokens have different meanings in different tasks.
It is not a hidden-state label. The learner does not receive target transition
tables, intermediate execution states, or correct call routes.

The shared neural state is a row distribution $\alpha\in\mathbb{R}^{1\times D}$.
Let $E\in\mathbb{R}^{N\times D}$ and $R\in\mathbb{R}^{D\times N}$ denote
encoder and readout logits. At temperature $\tau>0$, encoding and decoding use
row softmaxes, denoted $\overline E$ and $\overline R$. A bank contains $K$
learned operations and a provided identity operation $P_0=I_D$. The latter is an
execution substrate primitive, not a learned discovery. For token $a$ in context
$c$, a learned network produces $S$ slot distributions
$q_{c,a,s}(k)$, $k\in\{0,\ldots,K\}$. Soft execution is
\begin{align}
 \alpha_0 &= \overline E_{x,:},\\
 T_{c,a} &= \prod_{s=1}^{S}\left(\sum_{k=0}^{K}
                    q_{c,a,s}(k)\,\overline P_k\right),\\
 \alpha_{\ell} &= \alpha_{\ell-1}T_{c,a_\ell},
 \qquad \widehat p(y\mid x,w,c)=(\alpha_L\overline R)_y.
 \label{eq:growth-soft-execution}
\end{align}
Products are ordered left to right in execution order because states are row
vectors. Hard execution selects the encoder state, slot calls, operation
matchings, and decoded label. For a destination permutation $\pi_k$,
$(P_k)_{ij}=\mathbf{1}\{j=\pi_k(i)\}$; consequently $e_iP_k=e_{\pi_k(i)}$.
The continuous monodromy realization executes the selected operation through
path lifting. Its mathematical role and later coupling to learning are
developed in the monodromy section.

The compact bank stores destination coordinates
$\xi\in\mathbb{R}^{K\times D}$ and generates matching scores
\begin{equation}
 W_{kij}=-8(j-\xi_{ki})^2.
 \label{eq:compact-destination-scores}
\end{equation}
Every permutation is representable by setting $\xi_{ki}=\pi_k(i)$.
The persistent bank therefore contains $KD$ learned scalars, although forming
and differentiating its dense scores still incurs $KD^2$ temporary work and
storage. The ordinary constrained control receives this identical compression;
it is not a topology-specific saving.

With token vocabulary size $|\mathcal A|$, $C$ allocated contexts, embedding
width $H$, and $S$ slots, the implemented shared model has
\begin{equation}
 P_t=2ND_t+(|\mathcal A|+1+C)H
       +S(K_t+1)(H+1)+K_tD_t+4.
 \label{eq:growth-parameter-count}
\end{equation}
The final four scalars parameterize the shared geometric deformation. They are
absent from the ordinary constrained control. Preallocated context embeddings
and every frozen parameter are included. The reported runs use
$|\mathcal A|=6$, $C=6$, $H=48$, and $S=2$.

\subsection{Learning operations from terminal continuation profiles}

For a fixed context, write $F_c(x,w)$ for observed terminal feedback. A history
$h=(x,p)$ has a finite continuation profile
\begin{equation}
 \sigma_{\mathcal S}(h)=\big(F_c(x,ps)\big)_{s\in\mathcal S}.
 \label{eq:terminal-profile}
\end{equation}
The implemented observation table uses the empty word and all one-token words
as both prefixes and suffixes. Its successor check also consults words of
length three. All such observations come from the supplied TRAIN examples.
Different profiles provide explicit state-separation witnesses; equal profiles
are a provisional finite-data equivalence, not a claim of equivalence under
every future continuation.

The formation routine partitions observed histories by their profiles, checks
that observed successors close on the resulting classes, and rejects conflicting
or nonreversible successors. If the number of separated classes exceeds $D$,
the workspace grows. An assignment then aligns these classes to the current
encoder/readout coordinates. Already learned operations are reused when their
actions agree on every constrained reachable state. Otherwise, a constrained
matching completes an executable permutation, which is checked against the
entire current TRAIN set before neural installation. Candidate search,
matching, full-data checks, state expansion, and installation are all part of
learning cost. This procedure does not learn from hidden teacher states, but
does exploit the deliberately extensive short-word coverage of this benchmark.

After installation, the learner executes real gradient updates of the neural
encoder, readout, coordinates, embeddings, and calling network. In the measured
streams, new-task validation already equals one immediately after formation and
installation, before these updates. Thus the first acquisition in these runs is
attributable to terminal-profile identification and installation, rather than
to the subsequent 600 Adam updates alone.

\subsection{Binding protection and neural--native consistency}

Protecting an operation permutation alone does not protect a task: changing
the input encoding or call route can destroy an old computation even if every
bank operation remains unchanged. The implementation therefore stores admitted
native computation references and checks a transition-closed relation between
their reachable states and the updated computation. Shared proposals are
accepted, backtracked, or rejected according to this binding-level contract.
The relation permits state expansion and alternative internal coordinates when
they preserve decoded execution. It is a conventional semantic protection
condition, distinct from continuous-loop homotopy protection.

The soft interface requires an additional maintenance step. At each task
boundary, the implemented consistency refresh sharpens the model's own current
encoder, decoder, operation classes, and both original call slots. It obtains
no new task labels and adds no parameters. It preserves the current native
graph while increasing the probability assigned to its execution path. A
sufficient path-mass condition is stated in
Section~\ref{sec:theory-reproducibility}. Frozen models preserve aliased frozen
endpoints during this refresh. The same maintenance is available to ordinary
constrained and modular controls.

\subsection{Historical six-stage acquisition measurements}

Each stream contains six tasks. The first three use unknown reversible
$N$-state operations. The next three require a hidden-memory workspace with
$2N$ recurrent states and only $N$ observable output labels. The target hidden
mode and transition rules are never provided to the learner. TRAIN contains
all initial states paired with every word of length at most three, followed by
random words of lengths four through eight: 8,192 examples per task for
$N=8,16$ and 16,384 for $N=32$. The larger pool is required to include the
$259N=8,288$ mandatory short examples at $N=32$.

Validation, IID, OOD, and length-16 evaluations each contain 512 examples per
task. IID and validation use lengths at most eight and can overlap the exhaustive
short training coverage; reported OOD uses lengths nine through twelve.
Length-16 examples are disjoint from TRAIN by length. These are finite held-out
computations, not evidence about arbitrary lengths. This historical prototype is
bounded to lengths at most 18 and does not supply the required length-32/64
publication measurements.

Define $c_i(t)=1$ when task $i$ meets both neural and deployed thresholds at
the end of stage $t$: IID accuracy at least $0.95$ and both OOD and length-16
accuracy at least $0.90$. Admission additionally requires validation accuracy at
least $0.95$. The strictly retained boundary count is
\begin{equation}
 Q_t=\sum_{i\leq t}\prod_{u=i}^{t}c_i(u).
 \label{eq:boundary-retained-count}
\end{equation}
We report $\Delta P_t/\Delta Q_t$ only when a positive new stable capability
is obtained without losing an already stable one. The first acquisition pays
the entire initial allocation, $P_{-1}=0$. This is a stage-boundary criterion,
not an unobserved guarantee between measurements. These capability thresholds
define a reported outcome, not an accuracy-based eligibility test for baselines.

\begin{table}[t]
\centering
\small
\begin{tabular}{lrrrrrrl}
\toprule
Scale & $P_0$ & $P_1$ & $P_2$ & $P_3$ & $P_4$ & $P_5$ & Completed seeds\\
\midrule
$N=8$  & 1596 & 2232 & 2868 & 3946 & 4630 & 5314 & 0, 1, 2\\
$N=16$ & 2036 & 2720 & 3404 & 5130 & 5910 & 6690 & 0, 1, 2\\
$N=32$ & 3684 & 4464 & 5244 & 9034 &10006 &10978 & 0\\
\bottomrule
\end{tabular}
\caption{Historical native/coordinate-renderer streams, retained as internal
ablations rather than Primary Prometheus benchmark results. For every completed
seed shown, $Q_t=1,2,3,4,5,6$ and every seen-task neural/deployed accuracy equals
one on each of the four measured splits at each recorded stage boundary.
This table preserves the locally audited historical snapshot; it does not
complete the five-final-seed protocol for the future frozen Primary.}
\label{tab:growth-parameters}
\end{table}

At $N=8$, seed 0, the final bank contains 38 learned operations and $D=16$;
the analogous final workspaces are $D=32$ and $64$ at the larger scales.
The $N=32$ runs use a factored realization backend that preserves the learned
finite operation semantics without extra trainable parameters. This engineering
change and the differing acquisition pools prevent interpreting the table as
a controlled timing-scaling law.

Boundary maintenance resolves a measured neural/native mismatch, rather than
proving that the neural interface never forgets. In the repaired $N=8$, seed-0
stream, the minimum recorded old-task soft accuracies in the successive later
stages are approximately $0.348,0.309,0.252,0.242,0.178$, before returning to
one at the maintained boundary. Recorded deployed execution remains correct.
The original unrefreshed version retained only one task under the strict
boundary-stability criterion; its failed trajectory is preserved separately.

The strongest matched finite-program control is decisive for attribution:
ordinary constrained compact learning attains the same six-task boundary
result at $N=8$ and $N=16$, seed 0, with respectively 5,310 and 6,686 parameters,
exactly four fewer than the topology model. A repaired frozen-backbone policy
and separate task modules also attain six retained tasks at $N=8$, with 6,384
and 8,664 parameters. Therefore these measurements support compact shared
acquisition and a useful consistency repair, but do not isolate a monodromy
advantage in parameter growth for this historical version. In particular,
5,314 is not a measured parameter count for the future fully fused Primary.
Primary acquisition, retention, parameter growth, and all resource comparisons
must be established by actual retraining after the new architecture freeze.

\paragraph{Evidence locations.}
The frozen seed-0 growth and learning-curve extract is
\path{.deep-work/2026-10-31T234403-0400-v2/probes/topology-architecture/parameter-boundary-growth-evidence.json}
and
\path{.deep-work/2026-10-31T234403-0400-v2/probes/topology-architecture/parameter-boundary-learning-curves.csv}.
This extract predates the additional seeds. Their original summaries, raw
updates, checkpoints, and synchronization hashes are under
\path{.deep-work/2026-10-31T234403-0400-v2/probes/topology-architecture/runs/},
including \path{parameter-n8-seed1-topology-full-consistent},
\path{parameter-n8-seed2-topology-full-consistent},
\path{parameter-n16-seed1-topology-full-consistent}, and
\path{parameter-n16-seed2-topology-full-consistent}.
The $N=32$, seed-0 lineage is preserved in the
\path{parameter-n32-seed0-topology-through0-consistent},
\path{parameter-n32-seed0-topology-resumed-through2-consistent},
\path{parameter-n32-seed0-topology-resumed-through3-consistent}, and
\path{parameter-n32-seed0-topology-resumed-consistent} run directories.

\section{Learning executable monodromy classes}
\label{sec:monodromy-principle}

The mechanism studied here separates two operations that are easily conflated:
changing the implementation of an existing computation, and acquiring a
different executable computation. A deformation that is certified to preserve
a based-loop homotopy class can support the former, but cannot by itself
accomplish the latter. We therefore develop a learning interface with
class-changing proposals for unadmitted operations and class-preserving
transactions for admitted operations. The current evidence establishes this
interface in a small supplied covering; it does not yet establish an independent
topological advantage over equally informed conventional optimization.

\subsection{Covering space, marked states, and executable composition}

For an integer $D\geq 2$, consider the supplied polynomial family
\begin{equation}
 F_D(z,u)=z^D-Dz-u,
 \qquad
 \mathcal C_D=\{-(D-1)\zeta:\zeta^{D-1}=1\},
 \qquad U_D=\mathbb C\setminus\mathcal C_D.
 \label{eq:cover-and-critical-set}
\end{equation}
Projection onto $u$ defines an unbranched $D$-sheet covering over $U_D$.
We fix the control basepoint $u_*=0$ and mark its $D$ distinct roots once.
Lifting a based loop $\gamma:[0,1]\to U_D$ transports each marked root to
another root in the same fiber. Its endpoint action is a permutation
$\Pi_\gamma=\rho_D([\gamma])$. For chronological concatenation, defined so
that $\gamma_1\star\gamma_2$ executes $\gamma_1$ first,
\begin{equation}
 \Pi_{\gamma_1\star\gamma_2}
   =\Pi_{\gamma_2}\circ\Pi_{\gamma_1}.
 \label{eq:chronological-composition}
\end{equation}
An internal state is a sheet index, not a freshly initialized root at every
call. Equation~\eqref{eq:chronological-composition} consequently specifies an
executable, generally noncommutative composition rule.

The executed class-acquisition diagnostic uses $D=3$. The roots at $u_*=0$
are labelled $0,1,2$ in the order $(-\sqrt{3},0,\sqrt{3})$, and the critical
controls are $-2$ and $2$. Its one-parameter family is the polygonal loop
\begin{equation}
 \gamma_a:0\longrightarrow a+i\longrightarrow a-i\longrightarrow0,
 \qquad
 \Pi_{\gamma_a}=
 \begin{cases}
 B=(12),&a<-2,\\
 I,&-2<a<2,\\
 A=(01),&a>2.
 \end{cases}
 \label{eq:cubic-triangle-chambers}
\end{equation}
The singular cases $a=\pm2$ are not executable. For example,
$\Pi_B(\Pi_A(0))=2$, whereas $\Pi_A(\Pi_B(0))=1$; the order of calls matters.
The permutation labels in \eqref{eq:cubic-triangle-chambers} describe the
independently checked outcome, rather than an operation table provided to the
learner. The numerical executor actually continues all three roots with a
safeguarded predictor/Newton corrector. Its full-fiber endpoint observations,
known covering, critical controls, and numerical costs are available to every
comparison arm. Thus the protocol combines terminal task labels with
full-fiber execution observations; it is not a single-sheet-only observation
protocol. The executed source and independent reconstruction are preserved in
\path{.deep-work/2026-10-31T234403-0400-v2/probes/topology-architecture/versions/code-parameter-loop-boundary-learning-v1/loop_boundary_learning.py}
and
\path{.deep-work/2026-10-31T234403-0400-v2/supervision/reviews/review-038-loop-boundary-learning.md}.

\subsection{What safe implementation deformation does and does not learn}

An earlier construction first identifies a permutation, compiles it into a
loop of distinct root configurations, and then changes its geometry using a
time-dependent ambient shear. Writing a point as $x+iy$, this shear is
\begin{equation}
 \Phi_{\omega,t}(x+iy)
 =x+i\left[y+\sin^2(\pi t/T)h_\omega(x)\right],
 \qquad h_\omega(x)=\sum_\ell\omega_\ell\tanh(v_\ell x+b_\ell),
 \label{eq:safe-ambient-shear}
\end{equation}
where $v_\ell,b_\ell$ are fixed. Every finite $\omega$ gives a global
diffeomorphism, with an explicit inverse obtained by subtracting the same
vertical displacement. At both execution endpoints it is the identity.
Interpolating $\omega$ therefore gives a collision-free homotopy of
\emph{configuration-space} loops with a fixed base configuration, preserving
their braid and endpoint permutation. This is a useful within-class learning
freedom, but it is not a mechanism for discovering a new permutation. It is
also a different construction from the control-plane covering in
\eqref{eq:cover-and-critical-set}; the two spaces must not be identified
without an additional map between them.

The new class-learning interface instead trains the raw vertices of a loop in
the fixed control plane. A controller $a_\theta(e)$ maps a public operation
feature $e$ to coordinates in a supplied affine polygon family. For an
unadmitted operation, a finite update $\theta\mapsto\theta^+$ may place the
two regular endpoint snapshots in different chambers. The parameter
interpolation between them may pass through a singular loop, but that
interpolation is not executed. Only the separately checked regular snapshots
are run. There is consequently no claim that a nonsingular based homotopy can
change monodromy. The covering, basepoint, critical set and sheet labels remain
fixed. Moreover, a change of homotopy class in the kernel of $\rho_D$ would
not count as a new executable operation.

\subsection{Terminal boundary feedback and neural credit}

An exact endpoint-semantic loss is locally constant in regular loop
coordinates. Ordinary within-chamber differentiation alone therefore provides
no direct derivative for acquiring a different endpoint action. Let $V$ be
raw loop coordinates drawn from a smooth location kernel centered at
$a_\theta$, and let $\ell(V;\theta)$ denote the terminal loss obtained by
executing that loop and the current neural bindings. Define
\begin{equation}
 J_\sigma(\theta)
 =\int \ell(v;\theta)p_\sigma(v-a_\theta)\,dv.
 \label{eq:coupled-smoothed-objective}
\end{equation}
For the current fixed-cover, endpoint-semantic objective, the coupled
derivative has two contributions:
\begin{equation}
 \nabla_\theta J_\sigma
 =\mathbb E\!\left[\left.\partial_\theta\ell(V;\theta)
                    \right|_{V\ \mathrm{fixed}}\right]
 +(D_\theta a_\theta)^\top g_{\mathrm{boundary}}.
 \label{eq:coupled-boundary-derivative}
\end{equation}
The first term trains the input encoding, readout, calls and their shared
backbone while holding the actually observed endpoint maps fixed. For a
coordinate slice with crossing events $b_r$ and terminal-loss jumps
$\Delta\ell_r=\ell(b_r^+;\theta)-\ell(b_r^-;\theta)$, the second term uses
\begin{equation}
 (g_{\mathrm{boundary}})_j
 =\mathbb E_{V_{-j}}\!\left[
       \sum_r\Delta\ell_r\,p_\sigma(b_r-a_{\theta,j})\right].
 \label{eq:boundary-jump-estimator}
\end{equation}
Safe representatives of the neighboring chambers are actually executed to
obtain these terminal losses. The implementation uses a normalized
raised-cosine kernel and charges event detection, representatives, terminal
callbacks and continuation. A coordinate moving one free vertex affects its
two adjacent segments; tied coordinates moving several vertices are charged
for every affected segment instead of silently inheriting the smaller bound.

Equations~\eqref{eq:coupled-boundary-derivative}--\eqref{eq:boundary-jump-estimator}
are applications of known differentiation of piecewise-smooth objectives,
not a claimed new calculus theorem. A path-work objective with smooth
within-chamber dependence on $v$ would require its corresponding continuous
derivative as well. The implemented boundary term changes how probability
mass is assigned to executable classes; it does not replace successful
discrete execution by a mixture of several operations.

The ideal smoothing distribution assigns mass arbitrarily close to a
discriminant. The numerical executor additionally requires positive finite
clearance. These are distinct domains. The initial score-function arm was
censored by a probe entering the exclusion band, rather than assigned a
fabricated failure loss. The later hybrid interface evaluates an ideal
sample's class at a paid safe representative when available, while still
requiring the deterministic candidate loop itself to be executable. This
scalar-family convention is exposed to the conventional controls and is not
assumed to solve the general high-dimensional feasibility problem.

\subsection{A sufficient protected transaction}

For an admitted polygonal loop $\gamma$, let
$\delta=\min_{t,c\in\mathcal C_D}|\gamma(t)-c|>0$. If its updated vertices
satisfy
\begin{equation}
 \max_k|v_k(\theta^+,e)-v_k(\theta,e)|<\delta,
 \label{eq:protected-vertex-transaction}
\end{equation}
with the same segmentation and fixed endpoints, linear interpolation gives a
based homotopy avoiding $\mathcal C_D$. Indeed, the displacement of every
point on every interpolated segment is bounded by the maximum vertex
displacement; its distance from the critical set remains positive. Hence its
endpoint permutation is unchanged. Preserving each admitted generator then
preserves every word of those generators under fixed input/output bindings.

This sufficient topological certificate is accompanied by separate checks of
the executor's coordinate box and numerical clearance. In the diagnostic the
minimum allowed control clearance is $.015$. Protection of a permutation
alone does not protect a changing encoder, router or decoder; the additional
binding and soft-execution contracts are specified in
Section~\ref{sec:prometheus-block}. None of these conditions implies hardware
noise protection or a quantum energy advantage from a CPU simulation.

\subsection{Executed class acquisition and the unresolved advantage}

The initial scalar diagnostic uses two real shared weights,
\begin{equation}
 a_{\mathrm{old}}=w_0,\qquad
 a_{\mathrm{new}}=w_0+w_1,\qquad (w_0,w_1)=(-3,4).
\end{equation}
The old operation $B$ is supplied as the starting mechanism and verified by
continuation; it is not credited as learned in an earlier stage. Four
terminal training cases require $A$, which is outside the admitted subgroup
$\{I,B\}$. Two accepted neural updates produce
\begin{center}
\begin{tabular}{c|rr|rr}
snapshot & $a_{\mathrm{old}}$ & $a_{\mathrm{new}}$
         & training error & old error\\\hline
initial & $-3$ & $1$ & $.75$ & $0$\\
first & $-2.625$ & $1.75$ & $.75$ & $0$\\
second & $-2.26427259$ & $2.47145482$ & $0$ & $0$
\end{tabular}
\end{center}
The old loop actually changes geometry while retaining $B$; the new loop
crosses its class boundary and acquires $A$. Four frozen held-out composition
cases and an actual checkpoint reload also have zero error. These are finite
length-at-most-three tests, not a scaling or long-term continual-learning
result. The raw execution is
\path{.deep-work/2026-10-31T234403-0400-v2/probes/topology-architecture/runs/parameter-loop-boundary-learning-v1/learn/result.json}.

The equally informed ordinary event optimizer produces the same weights,
parameter hashes and counted trajectory. Discrete chamber search succeeds
with one update, $35$ program executions and $37$ terminal callbacks, versus
$70$ executions and $72$ callbacks for the event arm. Those totals include
their disclosed calibration and audit work. A reexecution guard spends more
continuation work than the analytic certificate on this instance, but the
ordinary optimizer is also allowed that certificate. Thus the result closes
the earlier map-first learning gap, while the principal question remains open:
does monodromy organization improve acquisition, retention, verification or
parameter growth beyond the strongest ordinary constraint and search methods
under matched information and complete cost?

\section{A trainable loop-based Prometheus Block}
\label{sec:prometheus-block}

The purpose of a Prometheus Block is to join operation formation with
executable state updates, learned calls, composition and retention in one
trainable state. The present implementation is a concrete fixed-three-state
prototype of that interface. It removes the requirement to learn a
permutation matrix first and render a loop afterward, but has not yet become
a general computational architecture. We specify its actual computation and
its observed failure before describing the remaining interface requirements.

\subsection{Neural state and computation}

The current \path{HybridLoopBlock} stores a trainable input encoder, a
trainable readout, call heads and a shared loop controller. A public feature
$e$ is transformed by a shared matrix $H$; a loop head $r$ produces
\begin{equation}
 h_\theta(e)=He,\qquad a_\theta(e)=r^\top h_\theta(e),
 \qquad
 \gamma_{a_\theta(e)}:0\to a_\theta(e)+i\to a_\theta(e)-i\to0.
 \label{eq:block-loop-controller}
\end{equation}
The two operation features and two token features are the disclosed vectors
$(1,0)$ and $(1,1)$. The same $H$ participates in both loop generation and
calling, so changing a new operation can change the old operation's geometry
and the call probabilities. For call slot $s$ and input token $b$,
\begin{equation}
 p_{\theta,s}(k\mid b)=\operatorname{softmax}(C_s h_\theta(e_b))_k,
 \qquad k\in\{\mathrm{skip},0,1\}.
 \label{eq:block-call-head}
\end{equation}
Skip supplies an identity computation. The other two choices refer to loop
bank entries. There is no trainable permutation-score tensor $W_k$ in this
prototype: endpoint maps are derived from paid actual root continuation.

Let $E_\theta$ and $R_\theta$ be row-stochastic matrices obtained from the
encoder and readout logits. For an input $x$ and a token word
$b_1\cdots b_T$, the soft learning path is
\begin{align}
 q_0&=E_\theta[x,:],\\
 q_{t,s+1}
   &=q_{t,s}\left[\sum_k p_{\theta,s}(k\mid b_t)M_k(V)\right],
       \qquad M_{\mathrm{skip}}=I,\\
 p_\theta(y\mid x,b_{1:T},V)&=[q_{T,S}R_\theta]_y,
 \label{eq:block-soft-execution}
\end{align}
where $M_k(V)$ is the one-hot matrix of the actually observed endpoint
permutation for the sampled loop coordinates. These matrices are held
constant during the within-chamber derivative in
Equation~\eqref{eq:coupled-boundary-derivative}. The terminal training loss is
$-\log p_\theta(y)$, with a declared numerical probability floor of
$10^{-12}$. The target supplies only the final task output, not a correct
route or an intermediate state.

Native execution uses the encoder, call and readout argmax choices and
composes the resulting permutations on the same sheet. This is the deployed
discrete program against which task success is checked. Cached maps are
derived observations of actual loops, not separately learned operations.
Reusing a receipt implements the same endpoint semantics without claiming
that the changed continuous trajectory has been freshly integrated. Cache
namespace, payload, root observations and cold-reexecution audits are
recorded and made available to conventional controls.

The executed one-slot model has $30$ trainable scalars:
$9$ encoder logits, $9$ readout logits, $4$ shared-backbone weights,
$2$ loop-head weights and $6$ call-head weights. It also stores $8$ fixed
feature-buffer scalars. Both the old operation and its call are supplied
initial conditions, while the new call begins uniformly across all three
choices. Encoder/readout logits begin near the identity but remain
trainable. This starting ability is disclosed rather than counted as
previously acquired learning. The exact executed source is
\path{.deep-work/2026-10-31T234403-0400-v2/probes/topology-architecture/versions/code-parameter-hybrid-loop-learning-v1/hybrid_loop_block.py}.

\subsection{One learning transaction, three retention obligations}

The persistent state comprises neural parameters $\theta$, admitted loop
receipts, a reference native binding relation, bounded derived-map storage
and the numerical execution contract. Given current terminal examples, a
learning transaction proceeds as follows.
\begin{enumerate}
 \item Sample the unadmitted operation's loop coordinates. Execute or reuse
       legitimately certified maps, charge their observations, and evaluate
       the actual soft terminal computation.
 \item Compute within-chamber gradients for encoding, readout and calls,
       together with the class-boundary term. Both derivatives pass through
       the same neural backbone; no correct new call is inserted manually.
 \item Propose an actual shared-weight update. Backtrack its finite scale
       while checking numerical executability, admitted loop protection,
       old native binding closure and soft/native agreement. Rejected
       executions and checks remain in the cost ledger. An aborted
       transaction restores the original parameter values in place.
 \item Admit a new operation and its token binding only after deterministic
       mean-loop execution gives perfect native and soft accuracy on the
       stated training admission set and satisfies the expanded confidence
       contract. The reference then records the actually acquired program,
       rather than a teacher operation table. Generalization remains a
       separate held-out test.
\end{enumerate}

These checks discharge different obligations. First, geometric clearance
protects the endpoint action of an old loop, as in
Equation~\eqref{eq:protected-vertex-transaction}. Second, an old task may
still be forgotten if its encoder or call changes, even when every bank
permutation is unchanged. The implementation therefore computes a finite
relation $\mathcal R$ between the reference and proposed native states,
starting from all old encoded input pairs, and requires
\begin{equation}
 \begin{split}
 &(E^-(x),E^+(x))\in\mathcal R,\\
 &(s,t)\in\mathcal R\ \Longrightarrow\
   R^-(s)=R^+(t),\quad
   (T_b^-(s),T_b^+(t))\in\mathcal R\quad(b\in\mathcal A_{\mathrm{old}}).
 \end{split}
 \label{eq:block-binding-closure}
\end{equation}
Here $E^\pm,R^\pm,T_b^\pm$ denote the complete hard native encoder,
readout and token transition before and after the update, including all call
slots. Exhausting this finite product-state closure preserves the decoded
behavior of every finite old word. It does not require the two internal
representations to use identical state labels. This is a conventional finite
behavioral-relation check; its correctness is not a topology-exclusive claim.

Third, native preservation alone is insufficient for the soft neural path.
Let $\underline p_E$, $\underline p_R$ and $\underline p_C$ be the minimum
probabilities assigned to the model's own hard encoder, readout and old-token
call choices. Because every executed bank map here is exactly one-hot, the
probability of following the entire own-native path through a word of length
at most $L$ is bounded below by
\begin{equation}
 m_L=\underline p_E\,\underline p_R\,
                 \underline p_C^{SL}.
 \label{eq:block-soft-native-mass}
\end{equation}
If $m_L>1/2$, the soft output has the same unique argmax as that native
program: the designated path alone carries more than half of the total
probability. The current transaction uses a small additional tolerance above
$1/2$ and $L=16$. The mathematical argument is real-arithmetic; the reported
float64 evaluation is not interval certification. Native closure covers all
finite old words, whereas this particular soft bound covers the specified
finite length. Keeping these scopes distinct prevents a finite confidence
check from becoming an unsupported unlimited-length neural guarantee.

\subsection{Composition and a constructive module interface}

The executable interface is a state transformation plus a receipt:
\begin{equation}
 (s,b,\theta,\mathcal M,\mathcal K)
 \longmapsto (s',y,\mathrm{receipt}),
 \qquad
 s'=\Pi_{k_S}\circ\cdots\circ\Pi_{k_1}(s),
 \label{eq:block-interface}
\end{equation}
where $\mathcal M$ is the admitted mechanism state and $\mathcal K$ contains
the execution and protection contracts. Consecutive calls reuse the output
sheet, making $AB$, $BA$ and longer words genuine computations rather than
independent classifications. Learning can revise the shared controller and
the call structure; admission converts a successfully acquired computation
into a protected reusable entry. Stored receipts permit certified semantic
reuse, with their maintenance and storage charged separately from parameters.

This interface also identifies the obligations for a general Block. A layer
must expose which state representation it passes to the next layer, and
learning must preserve the composed input--operation--readout binding, not
just each isolated loop. A capacity change must provide a state embedding,
transferred operation semantics and an updated protection relation. An
irreversible operation would require a disclosed reversible extension or a
different state update, charging the extra workspace, cleanup and readout.
Continuous inputs require a learned encoding and an explicit domain on which
its binding is protected. These are implementable interface requirements;
they are not achieved merely by calling the present $D=3$ module a general
Block. The prototype currently supplies fixed finite input features and a
fixed covering, and has not implemented capacity growth, noninvertible
computation, continuous-domain generalization or verified multilayer stacking.

\subsection{Executed hybrid V1: operation formation without task acquisition}

Hybrid V1 is an actual negative architecture-development result. The frozen
experiment uses four terminal training cases, four held-out composition
cases, a supplied old operation, and $100$ parameter transactions per arm.
The new call is initially uniform; success cannot be obtained by receiving a
correct route. The three arms apply the coupled gradient, only its
within-chamber component, or only its boundary component. Both components
are nevertheless computed and charged in all arms, so their run times are
not an ablation efficiency comparison.

\begin{center}
\begin{tabular}{l|rrrrr}
arm & updates & no-op & map
    & TRAIN (S/N) & held-out (S/N)\\\hline
coupled & $100$ & $72$ & $A$ & $.25/.25$ & $0/0$\\
within only & $100$ & $0$ & $I$ & $.25/.25$ & $0/0$\\
boundary only & $100$ & $94$ & $A$ & $.25/.25$ & $0/0$
\end{tabular}
\end{center}

All arms retain old-task soft and native accuracy $1$, and none admits the
new operation--binding pair. Acquiring the endpoint map $A$ therefore does
not establish acquisition of the task that should call it. A ``transaction''
is explicitly a committed proposal, including a zero-scale proposal; it must
not be reported as $100$ nonzero learning updates. Raw curves and actual
neural checkpoint reloads are preserved in
\path{.deep-work/2026-10-31T234403-0400-v2/probes/topology-architecture/runs/parameter-hybrid-loop-learning-v1/learn/result.json}.

The coupled arm exposes a specific repair target. At its final snapshot the
within-chamber call-head gradient norm is approximately $1.2583$, but every
tested positive joint step scale, down to $1/64$, is rejected by the old-loop
geometry/map guard; the transaction commits scale zero. Thus learning credit
exists while this single-direction transaction cannot use it. The observation
does not prove that every feasible update is blocked. It motivates testing a
component-wise or constrained repair that can preserve an admitted class
without discarding usable binding credit. Such a fallback is development
work, not an executed V1 result, and no outcome is predicted here.

The implementation accounts for terminal callbacks, root continuation,
event geometry, neural vector--Jacobian products, proposals and rejections,
finite relation checks, cache bytes and real checkpoint storage. Ordinary
controls receive the same mathematical certificates and execution
observations. Consequently neither the small parameter count nor old-task
retention in this failed acquisition test is evidence of an independent
Monodromy advantage. The concrete contribution at this stage is the unified
trainable interface and an observed, localized failure of its protection--
plasticity transaction, with the general-Block and competitive-benefit
questions still open.


\subsection{Executed hybrid V2: a targeted feasible-direction repair}

V2 retains V1's loss, examples, thirty trainable scalars, eight buffer scalars,
and protection certificates. Its only transaction change is to try the same
gradient on the existing encoder, readout, and call heads when all positive
joint backtracking scales are rejected. The backbone and loop head stay fixed
in that fallback; every proposed update still passes all three certificates.
This is a disclosed block-coordinate repair, applicable to ordinary constrained
optimization as well as this topological realization.

The complete arm finishes after $31$ gradient transactions and one own-class
confidence-maintenance transaction. Three commits use the head-only fallback.
TRAIN and the four held-out composition cases attain soft/native accuracy $1$;
old observed cases remain at $1$. Both single-component arms exhaust $100$
transactions without acquiring the new calling behavior. The first $28$
complete-arm transactions reproduce V1 exactly. New class acquisition occurs
before the eventual calling success, distinguishing the two learning events.

The complete arm uses $122$ executed programs, $716$ terminal callbacks,
$11{,}526$ continuation steps, and $127$ proposals, of which $95$ are rejected.
The one maintenance step sharpens its own acquired hard route without adding
parameters or receiving a teacher route. Its length-16 designated-path mass
rises from approximately $1.53\times10^{-5}$ to $.999779$. This validates the
specified interface on the diagnostic, not an architecture-level advantage or
a general multi-stage learner. Both gradient components are computed and charged
in every arm; their measured times do not rank optimized ablation efficiency.

The complete neural checkpoint and cold operation maps were actually reloaded.
Independent review039 checked all $67$ synchronized source/result files,
model tensors, raw curves, and the common prefix. The frozen run is
\path{probes/topology-architecture/runs/parameter-hybrid-loop-learning-v2/}
within the V2 campaign. This development component is not a frozen Primary
Prometheus and supplies no primary five-layer benchmark numbers.

% Roadmap correspondence: TODO 4. No real-data execution is claimed here.
\section{TODO 4: Transfer to Real-World Continual Learning}
\label{sec:real-world-transfer}

The same core computation and learning rule must be evaluated on algorithmic,
audio, and visual tasks. Input adapters and output heads may change, but changes
to the core block or its protection rule define an architecture variant and must
be reported separately. Synthetic permutation accuracy alone does not establish
this transfer.

\subsection{Fixed audio protocol}
Speech Commands v2 uses all 35 word classes in seven class-incremental stages,
with five new classes at each stage. Training, validation, and test examples
follow the official split and all methods use the same $16\,\mathrm{kHz}$ input
preprocessing. Each stage trains only on its permitted current data and declared
memory. Evaluation recognizes all classes seen so far without a task or class-group
identifier at inference.

\subsection{Fixed visual protocol}
Split CIFAR-100 uses ten class-incremental stages, each introducing ten fine-grained
classes. The official 50,000-example training set supplies a per-class 10\%
validation subset under a fixed seed; the official 10,000-example test set is not
used for development. All methods share preprocessing and augmentation rules and
use a single expanding classifier without a task identifier at inference.

\subsection{Memory, initialization, and cost}
Both tasks require Memory-0 and Memory-500 comparisons. The latter permits at most
500 old examples across all classes. Saved features, labels, teacher outputs,
statistics, and other persistent information are disclosed and charged separately;
a representation cache is not free memory. Main experiments use random
initialization. Pretrained runs, if any, form a separate comparison with matched
pretraining access. The fixed per-stage epoch budgets are 10, 50, and 200, with
validation-only configuration selection under the current benchmarking protocol.

\subsection{Evidence boundary}
No frozen-Primary audio or visual continual-learning result under these full
protocols is available. Earlier limited real-data development diagnostics are
retained separately and are not pooled into these formal comparisons.
In particular, the existing synthetic learner's public task embeddings cannot be
passed to these class-incremental evaluations as inference-time task identifiers.
The transfer study must test whether the learned computation actually reduces
acquisition, retention, parameter, or full lifecycle costs under these protocols.

\section{Benchmarking and competitive evaluation}
\label{sec:benchmarking}
\label{app:benchmarking}

\paragraph{Primary architecture and historical evidence.}
Primary Prometheus is the latest fully fused architecture actually validated
under TODO~3, with TODO~2's independent benefit also validated before the
Primary freeze. Its core mechanism, code version, and evaluation configuration
are frozen before actual retraining across all five benchmark layers. The
earlier native/coordinate-renderer streams are historical results and internal
ablations only; their outcomes, parameters, and costs are not pooled with or
substituted for new Primary measurements. The limited D3 hybrid V2 diagnostic
is development evidence, not completed TODO~3 validation or a Primary freeze.

\paragraph{Baseline validity is independent of accuracy.}
A configuration is valid evidence when its implementation is checked, any
claimed reproduction is faithful, information is fair, and its prospective
training/tuning protocol is followed and disclosed. No accuracy threshold
determines validity: 100\%, 37\%, and random-level accuracy are all retained
outcomes of valid configurations. Low performance can support a conclusion
conditional on the tested model, information, and budget; it does not establish
that a complete algorithm family has failed. An equal-quality parameter
comparison is simply unavailable when the measured qualities differ, rather
than making either run invalid. All failures, configuration trials, retries,
maintenance, and extra training remain charged. A low result is not a software
defect and does not trigger trials until a desired number appears.

\subsection{Five-layer evaluation scope}

\begin{table}[t]
\centering
\small
\begin{tabular}{p{0.07\linewidth}p{0.40\linewidth}p{0.43\linewidth}}
\toprule
Layer & Required scientific coverage & Current evidence boundary\\
\midrule
I & Six-stage acquisition at $N=8,16,32$ & Historical streams in
Table~\ref{tab:growth-parameters}; the frozen fully fused Primary must be
retrained under the five-seed protocol.\\
II & Noncommuting programs, unseen combinations, lengths 16/32/64, dynamic
write/overwrite/read/delay memory & Finite held-out words and length 16 tested;
the complete complex-program and dynamic-memory suite is not completed.\\
III & One Block for reversible, nonreversible, continuous, and variable-length
computation & Current finite reversible operations do not establish this
generality.\\
IV & Speech Commands v2 and Split CIFAR-100 & No competitive real-data result
is claimed in this manuscript checkpoint.\\
V & Applicable architectures, continual-learning controls, and mechanism
ablations & Historical valid measured configurations and working program
controls; prospective bounded family comparisons remain incomplete.\\
\bottomrule
\end{tabular}
\caption{All five layers require direct measurements. Evidence at one layer
does not replace another layer.}
\label{tab:benchmark-layers}
\end{table}

Layer II will distinguish operation-order dependence, including $AB\ne BA$
and commutators, from memorization of an observed word distribution. Programs
are trained at length at most eight and evaluated separately at lengths
16, 32, and 64. A separate key--value stream requires writing, overwriting,
reading, and delayed queries. Every architecture receives terminal task feedback
without teacher intermediate execution traces. Layer III keeps the Block's
core state-transition and learning mechanism fixed across task types; a change
to that core is reported as an architecture variant, rather than silently
counted as generalization.

\subsection{Information, development budgets, and final evaluation}

The synthetic comparison uses the same terminal data, stage order, public
context, and evaluation protocol. Dedicated constructors, teacher operator
tables, hidden state labels, and correct routes are not free priors for any
method. Formation passes and terminal-label accesses beyond minibatch training
are charged explicitly. Equal gradient-update counts do not imply equal total
learning compute when one method additionally solves a full observation table.

The locked synthetic update limits are B1: 600, B2: 2,400, and B3: 9,600
gradient updates per stage, available to every method including Prometheus.
Each main model family has at most 24 development configurations: development
seed 0 selects candidates, the top three are checked with development seed 1,
and the ranking uses a frozen validation criterion. After configuration lock,
the final evaluation uses seeds $0,1,2,3,4$ without test-based tuning. Shared
early checkpoints and resumed run fragments count as one lineage, not multiple
independent seeds. The real-data budgets are 10, 50, and 200 epochs per phase
under the same validation-only selection principle.

Only a localized, reproducible software, data, or training-logic error qualifies
as a bug fix. Optimizer, architecture, curriculum, or initialization changes
are configuration development unless they repair such a demonstrated error.
Any supplemental trajectory outside the registered main budget is disclosed
separately. Reaching the highest budget without acquisition remains an allowed
negative result. Test and OOD measurements do not select configurations,
training duration, or stopping points.

Three comparisons are kept distinct: identical information/protocol, matched
training and total-compute budgets, and equal attained capability quality.
The final analysis requires accuracy--parameter, accuracy--training-cost, and
accuracy--inference-cost frontiers. Different method families may use the
optimizers and schedules appropriate to their original algorithms; forcing one
recipe on all competitors is not a fairness requirement.

\subsection{Fixed real-data protocols and strong competitors}

Speech Commands v2 uses 35 word classes in seven class-incremental phases of
five classes. It uses the official train/validation/test split and common
16-kHz preprocessing. Each model classifies all seen words at inference
without a phase identifier or true class group. Split CIFAR-100 uses ten phases
of ten fine classes, the official 50,000 training and 10,000 test images, and a
fixed class-stratified 10\% training holdout for validation. It uses an expanding
single head and supplies no task ID at inference. Input processing,
augmentation, and data access are matched across architectures.

Both datasets have Memory-0 and Memory-500 conditions, where 500 is the total
number of stored old examples across all classes. Stored features, teacher
outputs, task statistics, and other historical state are additionally recorded;
they are not an exemption from the memory comparison. Main runs start from
random initialization. Pretrained variants are supplemental and give every
method the same pretraining information. These no-ID class-incremental
protocols are distinct from the public-context synthetic benchmark.

The primary competing architecture family is the Transformer: vanilla, replay,
adapters, and replay plus adapters, at two or more prospectively selected scales.
Parameter-matched and quality-matched comparisons are both required.
Continual-learning comparisons include DEN, PackNet, replay, and ER-ACE or an
equally mature class-incremental method. PackNet task-incremental results, when
applicable, are explicitly separated from no-ID class-incremental rankings.
Mamba-2 participates in appropriate sequence/memory tasks, and Titans in tasks
suited to its memory interface. HOPE/Nested Learning participates only with a
verifiable, applicable implementation. An unfaithful reproduction or untested
implementation is marked \texttt{NOT\_VALIDATED}, with its specific fidelity
or implementation gap stated. Low accuracy alone never assigns that status.

\subsection{Measured baseline results at the historical B1 budget}

\begin{table}[t]
\centering
\small
\begin{tabular}{lrrrrrrr}
\toprule
 & & \multicolumn{3}{c}{First task} & \multicolumn{3}{c}{New task at phase 6}\\
Method & Total $P$ & IID & OOD & L16 & IID & OOD & L16\\
\midrule
Transformer-64 &103752&.1836&.1230&.1348&.1152&.1094&.1016\\
Transformer-64+Replay &103752&.1836&.1230&.1348&.1191&.1172&.1055\\
Transformer-96 &229352&.1348&.1094&.1035&.1328&.1406&.1191\\
Transformer-96+Replay &229352&.1348&.1094&.1035&.1133&.1289&.1152\\
\bottomrule
\end{tabular}
\caption{Actual $N=8$, seed-0 historical streams, each with 600 updates per
phase and 3,600 updates overall. ``Phase 6'' is the newly introduced task at
the final phase, not an average over old tasks. These are valid measured B1
configurations under their reported historical protocol, regardless of their
low acquisition. They are not pooled with a future frozen Primary evaluation.}
\label{tab:historical-transformers}
\end{table}

All four runs in Table~\ref{tab:historical-transformers} have zero tasks passing
the dual-interface capability criterion. Their low first-task performance
shows that the result is not explained by forgetting alone. The $N=16$
width-64 Transformer similarly has no passing task in the measured stream.
These are valid budget- and configuration-conditional observations. They do
not by themselves characterize every Transformer configuration or establish an
equal-quality parameter ratio, and they are not new Primary Prometheus results.

A separately recorded anchored Transformer trial uses a 100-update source
prefix and 2,200 additional updates, totaling 2,300 updates. It is a supplemental
configuration trajectory, not a software bug fix or an equal-600-budget result.
It is valid configuration evidence under its own disclosed supplemental protocol.
The fixed development progression stops at its prospective budget. The
validation-selected frozen checkpoint is then tested once, obtaining IID
$0.37109375$, OOD $0.134765625$, and length-16 $0.111328125$. This improves
some short-word behavior without establishing full-task competence; the low
test result is not a trigger for unbounded additional trials.

The first 600-update task runs of the adapted DEN-RNN and linear
self-modifying CMS reconstruction obtain IID $0.115234375$ and
$0.138671875$, respectively. Their parameter counts are 1,800 and 81,224.
The former adapts DEN-style selection, expansion, splitting, and timestamps to
a recurrent terminal-feedback interface; the latter is not an official HOPE
reproduction. These are valid executed adaptations under their reported
600-update protocols; their low scores do not invalidate them. Conclusions
apply to these particular DEN-RNN and CMS variants. Faithful full-original
DEN and HOPE reproduction remains \texttt{NOT\_VALIDATED} because that fidelity
has not been established, not because an accuracy threshold was missed.

For the historical finite-program question, matched conventional controls are available.
At $N=8$, seed 0, ordinary constrained compact learning matches all six
stage-boundary capabilities using 5,310 parameters, versus 5,314 for the
topology model. Repaired freezing uses 6,384 and separate task modules 8,664.
The frozen repair adds a compatible 128-parameter readout when terminal-profile
output multiplicities require it, and later tasks reuse that readout. Its
initial failed version remains recorded. These controls explain the historical
shared/compressed parameter benefit without attributing it to monodromy; they
complement rather than replace the required Transformer comparison.

\subsection{Costs and causal attribution}

\begin{table}[t]
\centering
\small
\begin{tabular}{lrrrrr}
\toprule
 & Weight & Extra-state & Optimizer & Checkpoint & Old guard\\
Policy & bytes & JSON bytes & tensor bytes & bytes & seconds\\
\midrule
Topology compact &21256&208326&42544&341373&47.4276\\
Ordinary constraints &21240&104105&42508&202737&45.6787\\
Frozen + readout &25536&83598&7568&146821&0\\
Private modules &34656&81825&12316&160265&0\\
\bottomrule
\end{tabular}
\caption{Historical final-phase $N=8$, seed-0 resource fields for the
native/coordinate-renderer policies, not the future fully fused Primary.
The guard covers the phase's 600 updates.
These are format-specific fields and overlapping cost scopes, not a disjoint
sum or a controlled wall-time ranking.}
\label{tab:resource-fields}
\end{table}

All stored and frozen parameters count. Reference arrays are contained within
extra-state metadata; cache maps and certificates can also share storage.
Checkpoint sizes include optimizer, RNG, and container state. JSON size is not
minimal sufficient information or peak RAM. Phase wall time includes nested
training, evaluation, and protection scopes, which are not added twice.
Some jobs overlap on the GPU, so their individual wall times do not support a
throughput ranking. One ordinary-control initial-stage report is missing
post-training timing fields after a reload-helper failure; these remain null,
rather than being reconstructed as measured values.

Complete comparison additionally requires development and failed-run costs,
primitive discovery, data accesses, solver and installation work, temporary
dense tensors, replay, cache preparation, old-binding checks, rollback, and
checkpoint reload. Controlled same-hardware training FLOPs/GPU time,
end-to-end inference p50/p95, throughput, and peak memory are not yet complete.
Neither software simulation nor these parameter counts measure hardware energy
or establish a quantum advantage.

The causal matrix removes, separately, primitive reuse, sharing, workspace
growth, old capability protection, graph constraints, and native compilation.
Topology-specific comparisons substitute ordinary constraints for monodromy,
disable cross-class loop learning, disable within-class safe deformations,
and compare ordinary cross-class optimization and discrete search under the
same information and costs. Shared compression, program formation, graph
closure, and generic consistency repair are not credited to monodromy. A
same-mechanism reimplementation tests executable correspondence; it is not by
itself a historical-priority finding or a general impossibility result.

\subsection{Statistical reporting and completion status}

Formal comparisons use all five independent final seeds, mean, standard
deviation, and 95\% confidence intervals; paired-seed differences and their
intervals are primary. Multiple primary tests receive multiplicity correction,
and failed seeds remain included with their actual censoring or failure status.
The principal outcomes are final/average-incremental accuracy, forgetting,
backward transfer, acquisition cost, and total parameters and resources at
equal retained quality. The earlier scale/seed lineages remain historical or
internal ablations; they cannot count as retraining of the newly frozen Primary,
be pooled across architecture versions, or supply its five-seed statistics.

Publication-grade benchmarking remains \texttt{INCOMPLETE}: the fully fused
TODO~2 and TODO~3 must both pass; the complete architecture is then frozen
and actually retrained
across every layer against protocol-valid competitors under the bounded
development procedure. The two fixed real-data benchmarks, causal attribution,
and full resource/statistical evidence remain required. No competitor must
cross an accuracy gate to have its measured result included, and no favorable
historical finite-task result substitutes for the new Primary evaluation.

\paragraph{Evidence locations.}
Original neural streams and their actual checkpoints reside under
\path{.deep-work/2026-10-31T234403-0400-v2/probes/topology-architecture/runs/},
including \path{parameter-n8-seed0-transformer64-full},
\path{parameter-n8-seed0-transformer64-anchored-competence-v1},
\path{parameter-n8-seed0-den32-through0-v2}, and
\path{parameter-n8-seed0-selfmod_cms64-l1-through0-v2}.
The frozen stage-level extract is
\path{.deep-work/2026-10-31T234403-0400-v2/probes/topology-architecture/parameter-boundary-growth-evidence.csv}.
Actual original-source reading and adaptation boundaries for DEN are in
\path{.deep-work/2026-10-31T234403-0400-v2/sources/parameter-den/reading-and-adaptation.md}.
The governing prospective benchmark protocol is \path{ALL.md}, TODO~5, with
the latest explicit human directives controlling architecture identity and
accuracy-independent configuration validity if older wording conflicts.

\section{Theory, evidence boundaries, and reproducibility}
\label{sec:theory-reproducibility}
\label{app:evidence}

The role of theory is to identify the conditions under which an executable
operation, its task binding, and its soft neural realization remain correct.
These are distinct obligations. A mathematical invariant is not a proof of
task acquisition, a finite successful training stream is not a generalization
theorem, and an independent agent's agreement is not a scientific certificate.
The results below describe conditional finite mechanisms; they do not assert
an original general learning principle or a topology-specific cost advantage.
Historical native/coordinate-renderer results and limited hybrid diagnostics
remain separately versioned evidence. They are not measurements of Primary
Prometheus, which is defined by the latest fully fused TODO~3-validated
architecture and its subsequent frozen, actually retrained benchmark lineage.

\subsection{Continuation witnesses and finite-state identification}

\begin{proposition}[Observed state separation]
Fix a deterministic recurrent system and one task context. Suppose histories
$h$ and $h'$ have different observed outputs under the same continuation $s$.
Then any deterministic internal-state realization of both observations must
assign different recurrent states to $h$ and $h'$.
\end{proposition}
\begin{proof}
If their recurrent states were equal, applying the same continuation under the
same transition rule would produce equal final states and decoded outputs,
contrary to the observations.
\end{proof}

Thus $m$ pairwise distinct terminal continuation profiles require at least $m$
distinguishable states in this deterministic representation. This justifies
the historical prototype's witness-triggered workspace expansion. It does not show that equal
finite profiles are equivalent under all unobserved suffixes. The implemented
closure and full-TRAIN checks identify a model consistent with available
observations; its held-out behavior remains an empirical question. No
exponential global-closure conclusion is inferred from this finite witness.

\subsection{Binding-level protection under state expansion}

Let an old task model be $(X,\iota,T,r)$ and an updated model be
$(X',\iota',T',r')$, with a shared observable input domain and alphabet
$\mathcal A$. The encoders are $\iota,\iota'$, transitions are
$T_a,T'_a$, and readouts are $r,r'$. A relation
$\mathcal R\subseteq X\times X'$ protects the old computation if
\begin{align}
 &(\iota(x),\iota'(x))\in\mathcal R\quad\text{for every protected input }x,
 \label{eq:relation-encoding}\\
 &(z,z')\in\mathcal R\ \Longrightarrow\ r(z)=r'(z'),\\
 &(z,z')\in\mathcal R\ \Longrightarrow\
       (T_a(z),T'_a(z'))\in\mathcal R\quad(a\in\mathcal A).
 \label{eq:relation-closure}
\end{align}
\begin{proposition}[Native execution preservation]
Under these conditions, the old and updated systems have the same decoded
output for every finite word over the protected alphabet, starting from any
protected input.
\end{proposition}
\begin{proof}
The empty-word state pair belongs to $\mathcal R$ by
\eqref{eq:relation-encoding}. Each next token preserves membership by
\eqref{eq:relation-closure}. Induction on word length and readout agreement
give the result.
\end{proof}

The relation can connect different-sized state spaces, so it permits expansion
and internal reorganization without requiring every old parameter to freeze.
It protects equivalence to the admitted reference, not correctness of that
reference against an unknown task. Checking a finite reachable relation is
ordinary semantic verification, not by itself a topological contribution.
The reference construction, breadth-first exploration, acceptance/backtracking,
and stored relation all have measured maintenance costs.

\subsection{What continuous-loop protection establishes}

For the present permutation-valued realization, a loop in the unordered
configuration space of $D$ distinct points induces a permutation of the marked
fiber. With a fixed initial labeling, monodromy gives a representation
\begin{equation}
 \rho:\pi_1\!\left(\operatorname{Conf}_D(\mathbb C)/S_D\right)
       =B_D\longrightarrow S_D.
 \label{eq:permutation-monodromy}
\end{equation}
The quotient action records endpoint permutations; it is not a faithful
representation of the whole braid group. Different braid classes, including
pure braids, can have identical permutation actions. Ordinary graph equivalence
is likewise not a new topological invariant merely because it is used alongside
this realization.

\begin{proposition}[A clearance condition for based-loop retention]
Let $X=\mathbb C\setminus\mathcal B$ for a fixed finite branch set
$\mathcal B$ and a fixed covering over $X$. Let $\gamma$ be a polygonal based
loop with clearance
$\delta=\inf_{s,b\in\mathcal B}|\gamma(s)-b|>0$.
Let $\gamma'$ have the same parameterized polygon structure and the same
basepoint, and suppose the maximum corresponding vertex displacement is less
than $\delta$. Then $\gamma$ and $\gamma'$ are based-homotopic in $X$ and
induce the same marked-fiber monodromy.
\end{proposition}
\begin{proof}
Linear interpolation on each segment bounds every pointwise loop displacement
by the largest vertex displacement. For
$H(t,s)=(1-t)\gamma(s)+t\gamma'(s)$,
$|H(t,s)-b|\geq\delta-\sup_s|\gamma'(s)-\gamma(s)|>0$.
Thus the interpolation misses every branch point and fixes the endpoints.
Homotopy invariance of covering-space lifting gives the same monodromy.
\end{proof}

This is a conditional geometric protection rule: it permits real continuous
motion, rather than manually freezing all loop parameters. The covering,
singular set, basepoint, and labeling must satisfy the stated fixed conditions;
changing them requires its own argument. Applying it to a learned family also
requires a bound for every protected loop/input in the claimed domain.
Numerical continuation, finite-grid checks, and floating-point clearances are
not automatically rigorous interval certificates. A certified deformation
family may reuse an already paid-for realization across parameter changes, but
its assumptions, initial lift, invalidation, and checkpoint reconstruction must
all be included in the resource account.

Loop-class preservation, task-binding preservation, and accurate finite-time
execution remain separate. A gate can remain topologically unchanged while a
shared task router selects the wrong gate. Conversely, unchanged endpoint
behavior under a conventional constrained program is not evidence that its
learning benefit is topological. The software monodromy simulator is not a
physical quantum computer; no hardware energy or quantum speedup follows from
its numerical execution.

\subsection{A sufficient soft/native consistency bound}

Consider nonnegative stochastic soft execution as in
\eqref{eq:growth-soft-execution}. Suppose the hard computation uses at most
$S L_{\max}$ operation slots. Let $p_E,p_R,p_C,p_O$ be lower bounds on the
probabilities assigned, respectively, to the selected native encoding,
selected readout, every selected slot call, and every selected operation edge.
The probability mass following that particular native execution path is at
least
\begin{equation}
 p_{\mathrm{path}}\geq p_Ep_R(p_Cp_O)^{S L_{\max}}.
 \label{eq:soft-native-path-bound}
\end{equation}
If this bound exceeds $1/2$, the native output is the unique soft argmax for
every covered word, because all alternative output mass together is less than
$1/2$. The same conclusion can hold with a weaker bound, since several paths
may decode to the same output; this is only a sufficient condition.

The actual refresh evaluates these minima from the model's current soft
operation tensors and its own hard calls. It checks parameter count, original
slot structure, current native graphs, and frozen alias integrity. Subsequent
weight or temperature changes require renewed checking. Floating-point
evaluation of the analytical bound is not interval-certified verification.
The result transfers current native correctness to a bounded-length soft
interface; it does not teach a new class or prove always-online retention
between maintenance points. The rule is available to the non-topological
controls and therefore cannot establish a monodromy-exclusive effect.

\subsection{Executed evidence and version lineage}

Every reported run is tied to actual task arrays, configuration, code version,
raw update, formation, and guard records, curves, saved neural tensors, extra state,
and reload checks. Synchronization manifests supply file hashes for completed
copies. A successful exit alone does not demonstrate learning; measured
predictions and restored-checkpoint behavior supply the execution evidence.
Missing records are retained as missing. A completed server-side artifact and
a locally received runner exit receipt are distinct evidence: a transport
timeout does not authorize inventing the latter. Failed code, unsuccessful model
versions, and additional computation remain separate from the eventual
successful trajectory, rather than being relabeled as if success occurred
earlier.

The boundary evidence extraction reads original JSON/JSONL records and safe ZIP
metadata without training or executing checkpoint code. Its CSV exposes
stage-boundary outcomes and intermediate curves separately. Subsequent
additional seeds remain traceable to their own original run manifests rather
than being silently attributed to that older extract. Frozen-backbone storage
audits additionally check tensor aliasing and unchanged old weights; these
audits establish stored values and counts, not independent neural inference.
The raw GPU reload records provide the latter within the execution environment.

Inference reload is also distinct from exact training continuation. In the
limited hybrid V2 diagnostic, actual neural weights and cold reexecuted maps
were checked after reload, but the external Python smoothing RNG state was not
saved. That checkpoint therefore establishes inference restoration, not
bit-exact interrupted training recovery. Exact recovery requires the relevant
RNG state, sampler/transaction position, and any optimizer or pending-update
state, with an actual continuation check before making that stronger claim.

\paragraph{Reproducibility artifacts.}
The source, data generators, and raw run tree are under
\path{.deep-work/2026-10-31T234403-0400-v2/probes/topology-architecture/}.
The historical finite formation implementation is \path{terminal_profile_formation.py},
the historical compact model is \path{parameter_native_models.py}, the experiment driver
is \path{parameter_efficiency_study.py}, and neural/native maintenance is
\path{parameter_consistency_refresh.py} within that directory. Independent
within-campaign examinations are preserved at
\path{.deep-work/2026-10-31T234403-0400-v2/supervision/reviews/review-037-boundary-growth-and-frozen-readout.md}
and
\path{.deep-work/2026-10-31T234403-0400-v2/supervision/reviews/review-038-loop-boundary-learning.md}.
The latter concerns a distinct, small terminal-feedback loop-class diagnostic;
it does not substitute for the six-stage architecture or publication benchmark.
The subsequent integrated hybrid V1/V2 evidence and its inference-versus-RNG
recovery boundary are examined in
\path{.deep-work/2026-10-31T234403-0400-v2/supervision/reviews/review-039-hybrid-loop-development.md}.
That limited D3 development result is not full TODO~3 validation or a Primary
architecture freeze.

The frozen growth extract's SHA256 is
\path{311a5073e8f777fc6c6270d5323ebee258542b9be15a50d1bae2fa3ddb3e6c12},
and its learning-curve CSV has SHA256
\path{b7e9a6b9e11ce327e0e2b1d108ae3ea95b82e3f1483f32cc7e51d50f88af28b4}.
These identify particular evidence snapshots, not a claim that every later
artifact or current source has the same hash. Each executed version and later
run must retain its own provenance.

\subsection{Independent checking and limits of present claims}

The campaign uses independent agent contexts to inspect original equations,
code, raw outputs, costs, and saved tensor metadata. Such inspection has
corrected coding, physical-model, cost, and claim-scope errors. It is not
independent laboratory replication, community peer review, or proof by model
consensus. The finite propositions above are mathematical arguments in the
stated models; this manuscript does not claim Lean compilation or formal
verification of their implementations. Primary-source reading records preserve
what was actually read and the boundaries of adaptations. Historical novelty
requires a checked comparison with substantive original work, beyond an
independent reimplementation of the same rule.

The final reproducibility package requires fixed data splits, all five final
seeds, locked configurations, executable environments, checkpoints, curves,
raw failures, causal ablations, and statistical analysis scripts. Formal
confidence intervals and paired comparisons remain to be produced under the
protocol in Section~\ref{sec:benchmarking} after the fully fused Primary's core,
code, and configuration freeze. Historical architecture versions cannot be
pooled into that lineage or selected as its substitute. Baseline configuration
validity depends on implementation, claimed reproduction fidelity, fair
information, and the prospective protocol, never on attaining an accuracy
threshold. Low-scoring valid configurations and unsuccessful development
trials remain in the evidence. Complete cost reporting must include
development, discovery, formation, maintenance, storage, verification, and
inference in addition to gradient updates.

At this checkpoint, historical evidence supports finite executable continual
acquisition with boundary consistency and matched ordinary controls; the
separate hybrid diagnostic supports a limited integrated learning repair.
Neither supplies newly retrained Primary results. An independently useful
topological learning contribution, a general
nonreversible/continuous Block, competitive real-data results, and the complete
publication-grade evidence package remain unresolved. This manuscript is a
working scientific record, not a declaration of breakthrough, publication
readiness, or public release; authorship and intellectual-property conditions
must be settled before external dissemination.


\subsection{Training-state recovery: executed new audit}

The original hybrid V1/V2 checkpoints support inference reload but omit their
external smoothing RNG state. That missing history remains missing. A new
checkpoint wrapper additionally records the Python sampling generator, Torch
CPU RNG, tensor bits and strides, parameter gradient flags, training control,
source hashes, semantic receipts, cache state, and logical work counters.

A separately seeded audit starts from the actual V2 transaction-30 checkpoint,
performs one new gradient calculation, saves, and compares the next continuous
step against the next step after disk reload. The three actual new gradient
calculations match in all ten frozen comparison groups, including both credit
components, complete gradients, neural tensors, RNG, execution outputs, semantic
state, and work counts. Physical restore/constructor overhead is recorded
separately from restored logical training history. This is a CPU explicit-SGD
recovery check, not reconstructed original V2 training or a CUDA guarantee.

The first audit attempt failed before any gradient calculation because the
isolated process could not read a checkpoint placed beside its inspected source
files. The retry places the identical input in the runner's existing read-only
input allowance; isolation and source rules are unchanged. Both attempts and
their actual exit receipts are retained under
\path{probes/topology-architecture/runs/parameter-hybrid-training-recovery-v1/}
and \path{parameter-hybrid-training-recovery-v2/} in the same campaign.


\end{document}
